\documentclass[11pt,a4paper,openany]{book}

% ============================================================
% Packages
% ============================================================
\usepackage[margin=1in]{geometry}
\usepackage{amsmath,amssymb,amsthm}
\usepackage{mathtools}
\usepackage{booktabs}
\usepackage{array}
\usepackage{tabularx}
\usepackage{longtable}
\usepackage{graphicx}
\usepackage{float}
\usepackage{enumitem}
\usepackage{fancyhdr}
\usepackage{titlesec}
\usepackage{hyperref}
\usepackage[dvipsnames]{xcolor}
\usepackage{tcolorbox}
\usepackage{listings}
\usepackage{caption}
\usepackage{setspace}
\usepackage{microtype}

% ============================================================
% Colors & Styles
% ============================================================
\hypersetup{
    colorlinks=true,
    linkcolor=MidnightBlue,
    citecolor=OliveGreen,
    urlcolor=BrickRed
}

\definecolor{codegreen}{rgb}{0.25,0.49,0.31}
\definecolor{codegray}{rgb}{0.5,0.5,0.5}
\definecolor{codepurple}{rgb}{0.58,0,0.82}
\definecolor{backcolour}{rgb}{0.97,0.97,0.97}

\lstdefinestyle{pythonstyle}{
    backgroundcolor=\color{backcolour},
    commentstyle=\color{codegreen},
    keywordstyle=\color{codepurple}\bfseries,
    numberstyle=\tiny\color{codegray},
    stringstyle=\color{BrickRed},
    basicstyle=\ttfamily\small,
    breaklines=true,
    numbers=left,
    numbersep=5pt,
    frame=single,
    framesep=5pt,
    xleftmargin=15pt,
    language=Python,
    showstringspaces=false
}
\lstset{style=pythonstyle}

% ============================================================
% Theorem-like environments
% ============================================================
\theoremstyle{definition}
\newtheorem{definition}{Definition}[chapter]
\newtheorem{insight}[definition]{Key Insight}

\theoremstyle{remark}
\newtheorem*{takeaway}{Takeaway}

% ============================================================
% Colored boxes
% ============================================================
\newtcolorbox{keyidea}[1][]{
    colback=blue!5!white,
    colframe=MidnightBlue,
    title=#1,
    fonttitle=\bfseries,
    sharp corners,
    boxrule=0.5pt
}

\newtcolorbox{warning}[1][]{
    colback=red!5!white,
    colframe=BrickRed,
    title=#1,
    fonttitle=\bfseries,
    sharp corners,
    boxrule=0.5pt
}

% ============================================================
% Headers
% ============================================================
\pagestyle{fancy}
\fancyhf{}
\fancyhead[LE]{\leftmark}
\fancyhead[RO]{\rightmark}
\fancyfoot[C]{\thepage}
\renewcommand{\headrulewidth}{0.4pt}

% ============================================================
% Title
% ============================================================
\title{
    \vspace{-2cm}
    {\Huge\bfseries Truly Master\\Reinforcement Learning}\\[1em]
    {\Large\itshape From Bellman Equations to Preference Learning:\\
    A Conceptual Deep Dive}\\[2em]
    {\large Chen Z.}
}
\date{\today}
\author{}

% ============================================================
% Document
% ============================================================
\begin{document}

\frontmatter
\maketitle

% ------ Preface ------
\chapter*{Preface}
\addcontentsline{toc}{chapter}{Preface}

You have read Sutton \& Barto. You know what Q-learning is. You can recite the Bellman equation. You have maybe even implemented PPO.

But something still feels off. Why did RL move from value functions to policy gradients? What exactly is the relationship between DQN, PPO, SAC, and RLHF? Why do modern methods look so \emph{ad-hoc} compared to the elegant Bellman theory?

This book answers those questions. It is not another algorithm-by-algorithm tutorial. It is a \textbf{conceptual deep dive} that traces one unified story: how reinforcement learning evolved from classical dynamic programming to modern preference learning, and why every algorithm along the way is a natural consequence of that evolution.

We approach RL from two complementary angles:

\begin{itemize}
    \item \textbf{The macro view}: a single narrative arc from Bellman (1957) to GRPO (2025), showing how the Bellman equation went from ``the entire optimization'' to ``completely absent.''
    \item \textbf{The micro view}: a multi-dimensional dissection of \emph{why} each technique exists---loss design, sampling structure, bias management, compression, and action modeling.
\end{itemize}

The goal is not completeness---it is compression. We aim for the shortest path from ``I know what Q-learning is'' to ``I understand why the entire field evolved the way it did.''

\vspace{1em}
\noindent\textit{This book is for students who passed the RL course but can't connect the dots, engineers who want deeper intuition, and anyone preparing for interviews who is tired of surface-level algorithm summaries.}

\tableofcontents

% ============================================================
\mainmatter
% ============================================================

\part{The Big Picture}
\chapter{RL Is Not What You Think}
\label{ch:rl-not-what-you-think}

Most introductions to reinforcement learning begin with the Bellman equation, march through Q-learning, and end with a collection of algorithms.
This framing creates a dangerous illusion: that RL is merely supervised learning with delayed labels.
It is not.
The single most important thing to understand about RL is its \emph{sampling structure}---and how that structure differs from supervised learning at a fundamental level.

% ===========================================================
\section{The Same Optimization Framework}
% ===========================================================

At the level of code, RL and supervised learning (SL) are remarkably similar.
Both minimize a loss with gradient descent:
%
\begin{equation}
\theta \leftarrow \theta - \alpha\,\nabla_\theta \mathcal{L}(\theta).
\end{equation}
%
In SL the loss is typically a prediction error:
$\mathcal{L} = L\bigl(f_\theta(x),\, y\bigr)$.
In RL the loss takes one of two canonical forms:
%
\begin{itemize}
    \item \textbf{Value-based:}
        $\mathcal{L} = \bigl(Q_\theta(s,a) - y\bigr)^2$,
        where $y = r + \gamma \max_{a'} Q_{\bar\theta}(s',a')$.
    \item \textbf{Policy-based:}
        $\mathcal{L} = -\log\pi_\theta(a|s)\,A(s,a)$.
\end{itemize}
%
Both use SGD or Adam.  Both backpropagate through the same autodiff graph.
At the ``write a training loop'' level, RL and SL share the same optimization framework:
forward pass, compute loss, backward pass, update.

\begin{keyidea}[Surface Similarity]
RL and SL share the same \texttt{loss.backward(); optimizer.step()} pipeline.
The difficulty of RL is \emph{not} in the optimizer---it is in where the data comes from.
\end{keyidea}

% ===========================================================
\section{The Closed Loop: Where Everything Changes}
\label{sec:closed-loop}
% ===========================================================

In SL, the dataset $\mathcal{D} = \{(x_i, y_i)\}$ is fixed before training begins.
The model changes; the data does not.
This is an \textbf{open-loop} structure:
%
\begin{equation}
\text{Fixed } \mathcal{D} \;\longrightarrow\; \mathcal{L} \;\longrightarrow\; \nabla_\theta \;\longrightarrow\; \theta' \;\longrightarrow\; \text{better model.}
\end{equation}

In RL, data is generated by the policy itself.
When $\theta$ changes, the policy $\pi_\theta$ changes, which changes the sampling distribution, which changes the data, which changes the loss landscape.
This is a \textbf{closed loop}:
%
\begin{equation}
\pi_\theta \;\longrightarrow\; \text{sample data} \;\longrightarrow\; \mathcal{L} \;\longrightarrow\; \nabla_\theta \;\longrightarrow\; \pi_{\theta'} \;\longrightarrow\; \text{new data distribution.}
\end{equation}

\begin{keyidea}[The Closed Loop]
In SL, the algorithm does not affect the data distribution.
In RL, the algorithm \emph{is} the data distribution generator.
The policy is simultaneously the object being optimized and the source of training data.
\end{keyidea}

An analogy helps sharpen this:
\begin{itemize}
    \item \textbf{SL} is a student with a pre-printed textbook.
          How the student studies does not change the textbook.
    \item \textbf{RL} is a student who writes their own textbook as they learn.
          What they study affects what practice problems appear next,
          which affects what they learn after that.
          A premature bias produces a biased textbook for all future study.
\end{itemize}

% ===========================================================
\section{Four Links in the Causal Chain}
\label{sec:four-links}
% ===========================================================

The closed loop can be decomposed into four causal links, each feeding into the next:
%
\begin{enumerate}
    \item \textbf{Policy determines sampling.}
          The current $\pi_\theta(a|s)$ decides which state--action pairs are visited.
    \item \textbf{Sampling determines data.}
          The visited pairs form the empirical distribution on which all estimates are based.
    \item \textbf{Data determines value estimates.}
          $Q$, $V$, and advantage $A$ are all estimated from the sampled data---they inherit its biases and gaps.
    \item \textbf{Value estimates update the policy.}
          Whether through $\arg\max Q$ or advantage-weighted log-likelihood, the estimated values reshape $\pi_\theta$.
\end{enumerate}
%
Then the cycle repeats.
Every RL-specific technique---replay buffers, target networks, importance sampling, clipping, entropy bonuses---exists to manage pathologies that arise \emph{somewhere in this loop}.

\begin{table}[ht]
\centering
\caption{RL techniques mapped to the closed-loop problems they address.}
\label{tab:closed-loop-fixes}
\begin{tabular}{ll}
\toprule
\textbf{Technique} & \textbf{Closed-Loop Problem Addressed} \\
\midrule
Replay buffer       & Temporal correlation; distribution drift \\
Target network      & Bootstrap target chasing itself \\
$\epsilon$-greedy / entropy bonus & Sampling collapse; insufficient coverage \\
Importance sampling & Distribution mismatch (old vs.\ new policy) \\
PPO clipping        & Excessive policy change per step \\
Double Q-learning   & $\max$ operator amplifying estimation noise \\
\bottomrule
\end{tabular}
\end{table}

None of these techniques are necessary in SL, because SL's data distribution is not affected by the model.

% ===========================================================
\section{RL Labels Are Not Given---They Are Constructed}
% ===========================================================

In SL, labels $y$ are externally provided and independent of the model.
In RL, the ``labels'' are constructed by the learning system itself:
%
\begin{itemize}
    \item In value-based RL, the target $y = r + \gamma \max_{a'} Q_{\bar\theta}(s', a')$ depends on the model's own Q-estimates.
          The model creates its own labels and then fits them---\emph{self-bootstrapping regression}.
    \item In policy gradient methods, the advantage $A(s,a)$ is estimated from sampled returns or a learned critic, both of which depend on the current policy.
\end{itemize}
%
This self-referential structure means that errors compound: a biased Q-estimate produces biased targets, which produce more biased Q-estimates.
Understanding this recursive dependency is essential for understanding why RL is unstable in ways that SL is not.

% ===========================================================
\section{The Real Definition of RL}
\label{sec:real-definition}
% ===========================================================

With the closed loop in mind, we can state what RL truly is---not in textbook terms, but in operational terms:

\begin{keyidea}[RL in One Sentence]
Reinforcement learning is \textbf{biased preference optimization under shifting distributions}:
on a constantly changing data distribution, using biased sampled data,
estimate long-term returns, and optimize the next round's sampling strategy.
\end{keyidea}

This definition highlights three properties that separate RL from SL:
\begin{enumerate}
    \item \textbf{Shifting distributions.}  The data distribution is non-stationary because the policy generates it.
    \item \textbf{Biased estimates.}  Q-values, advantages, and returns are all estimates with sampling noise, bootstrap bias, and function approximation error.
    \item \textbf{Preference optimization.}  The goal is not to find a ``true'' value function in some god's-eye sense, but to stably push the behavioral distribution toward actions that are statistically better.
\end{enumerate}

This is why modern RL has drifted away from the Bellman equation as a primary optimization target and toward methods that directly shape the policy distribution (PPO, SAC, RLHF, GRPO).
The Bellman equation describes an ideal fixed point; the training process is a noisy, nonlinear, recursive iteration that must be stabilized by engineering.

% ===========================================================
\section{Why This Matters for Everything That Follows}
% ===========================================================

The rest of this book traces two consequences of the closed loop:

\begin{enumerate}
    \item \textbf{The macro arc} (Chapter~\ref{ch:grand-arc}): how the field evolved from dynamic programming---where the Bellman equation is everything---to DPO/GRPO, where the Bellman equation is completely absent.
          Each stage is a response to the closed loop's instabilities.
    \item \textbf{The micro mechanisms} (Part~II): how specific techniques (loss design, on/off-policy structure, exploration, bias control, compression, action modeling) each address a particular link in the four-step causal chain.
\end{enumerate}

The unifying theme is simple:
\emph{RL is not hard because the optimization is hard.
RL is hard because the ground shifts under the optimizer's feet.}
Every algorithm in this book is, at its core, a strategy for optimizing on shifting ground.

\chapter{From Bellman to Preference Learning: The Grand Arc}
\label{ch:grand-arc}

This chapter tells a single story.
It begins in 1957 with Richard Bellman's dynamic programming and ends in 2025 with DeepSeek's GRPO---a method in which the Bellman equation appears nowhere.
Along the way, we trace nine stages of evolution and identify the forces that drove each transition.

The thesis in one sentence:

\begin{keyidea}[The Grand Arc]
Classical DP solves for the global optimum under a known model.
Modern RL learns a good-enough, stable, preference-aligned policy under limited sampling, noise, bias, and distributional uncertainty.
These two objectives differ in nature: the former is god's-eye analytical optimization; the latter is statistical preference learning.
\end{keyidea}

% ===========================================================
\section{Stage 1: The Bellman Equation---Correct, but Ideal}
\label{sec:stage1-dp}
% ===========================================================

In 1957, Richard Bellman introduced dynamic programming and the principle of optimality.
In a deterministic, fully known environment with transition model $P(s'|s,a)$ and reward $r(s,a)$, the Bellman equation is both necessary and sufficient.

For a given policy $\pi$:
\begin{equation}
V^{\pi}(s) = \mathbb{E}_{a \sim \pi,\, s' \sim P}\bigl[r(s,a) + \gamma\, V^{\pi}(s')\bigr].
\end{equation}
For optimal control:
\begin{equation}
V^{\ast}(s) = \max_a\, \mathbb{E}_{s' \sim P}\bigl[r(s,a) + \gamma\, V^{\ast}(s')\bigr].
\end{equation}
Or in Q-function form:
\begin{equation}
\label{eq:bellman-q}
Q^{\ast}(s,a) = \mathbb{E}_{s' \sim P}\bigl[r(s,a) + \gamma\, \max_{a'} Q^{\ast}(s',a')\bigr].
\end{equation}

These equations are mathematically rigorous.
The Bellman equation \emph{is} the entire optimization---solving it yields the globally optimal policy.
The problem is not that the Bellman principle is wrong, but that it describes an \textbf{ideal fixed point}, while the training process is not a fixed point.

% ===========================================================
\section{Stage 2: Stochastic DP and MDPs (1960)}
\label{sec:stage2-mdp}
% ===========================================================

Ronald Howard extended Bellman's framework to stochastic environments (Markov Decision Processes) and introduced Policy Iteration and Value Iteration.
The essence did not change---still finding the Bellman equation's fixed point under a known model---but transition probability uncertainty was now accommodated.

The Bellman evaluation equation becomes:
\begin{equation}
V^{\pi}(s) = \sum_a \pi(a|s) \sum_{s'} P(s'|s,a)\bigl[r(s,a) + \gamma\, V^{\pi}(s')\bigr].
\end{equation}

This operator is fundamentally \emph{averaging}: it sums over actions weighted by the policy and over next states weighted by transition probabilities.
Averaging is smooth, stable, and easy to prove contraction for.
This stability is why V-functions were developed first.

% ===========================================================
\section{Stage 3: TD Learning---From Known Models to Sampling (1988)}
\label{sec:stage3-td}
% ===========================================================

Sutton's Temporal Difference (TD) learning was the first major departure from the god's-eye view.
The key move: replace exact expectations with sampled estimates.

\begin{equation}
V(s) \leftarrow V(s) + \alpha\bigl[r + \gamma\, V(s') - V(s)\bigr].
\end{equation}

No model is needed---only the ability to observe transitions $(s, a, r, s')$.
V-functions were the natural starting point because they average over action outcomes, yielding low-variance, stable updates.

The cost: V cannot directly tell you which action to take.
To extract a policy from $V$, you need the transition model $P(s'|s,a)$---precisely what you no longer have.

% ===========================================================
\section{Stage 4: Q-Learning---Making Action Optimization Explicit (1989)}
\label{sec:stage4-q}
% ===========================================================

Watkins' Q-learning solved the action-selection problem by shifting the learning target from $V(s)$ to $Q(s,a)$:
\begin{equation}
Q(s,a) \leftarrow r + \gamma\,\max_{a'} Q(s',a').
\end{equation}

Policy improvement becomes trivial: $a^{\ast} = \arg\max_a Q(s,a)$.
No environment model needed.
Q-learning is also \textbf{off-policy}---the update assumes the future will choose optimally regardless of what the behavior policy actually did.
This makes it data-efficient: a high-value signal can propagate rapidly.

\subsection{The Cost of $\max$}

But the $\max$ operator is an unbridled horse.
During training, Q is a noisy estimate: $\hat{Q} = Q^{\ast} + \epsilon$.
Therefore:
\begin{equation}
\max_a \hat{Q}(s,a) = \max_a \bigl[Q^{\ast}(s,a) + \epsilon(s,a)\bigr].
\end{equation}
The $\max$ preferentially selects the action with the largest positive noise.
This creates a feedback loop:
\begin{enumerate}
    \item An action is overestimated due to noise.
    \item $\max$ selects it for the Bellman target.
    \item The inflated target is written back into the Q-network.
    \item Future updates propagate the overestimation further.
\end{enumerate}

\begin{keyidea}[The Core Tension]
\textbf{Averaging} (V-style) is stable but conservative---it does not amplify the best action.
\textbf{Max} (Q-style) is data-efficient but aggressive---it amplifies estimation noise.
The entire history of RL oscillates between these two poles.
\end{keyidea}

\begin{table}[ht]
\centering
\caption{The averaging--max tension.}
\label{tab:avg-vs-max}
\begin{tabular}{lll}
\toprule
\textbf{Approach} & \textbf{Strengths} & \textbf{Weaknesses} \\
\midrule
Averaging / expectation & Stable, smooth, low variance & Slow learning, weak optimization \\
Max / greedy & Data-efficient, fast propagation & Overestimation, instability \\
\bottomrule
\end{tabular}
\end{table}

% ===========================================================
\section{Stage 5: DQN---Taming Qmax with Deep Networks (2013--2016)}
\label{sec:stage5-dqn}
% ===========================================================

DQN did not invent Q-learning.
Its true significance is the first reasonably stable combination of deep neural networks with off-policy Q-learning.
The Bellman target remained unchanged:
$y = r + \gamma\,\max_{a'} Q_{\bar\theta}(s',a')$.
DQN tamed the $\max$ operator through engineering:

\begin{enumerate}
    \item \textbf{CNN encoder}: map high-dimensional pixel states to features.
    \item \textbf{Replay buffer}: break temporal correlation, improve data reuse.
    \item \textbf{Target network} $\bar\theta$: prevent the target from chasing itself, reducing bootstrap instability.
    \item \textbf{Off-policy learning}: reuse past data to train the current value function.
\end{enumerate}

\subsection{Subsequent Repairs}

DQN exposed the $\max$ operator's weaknesses at scale.
Two key follow-ups:

\textbf{Double DQN} decouples action selection from evaluation:
\begin{equation}
a^{\ast} = \arg\max_a Q_\theta(s',a), \qquad
y = r + \gamma\, Q_{\bar\theta}(s', a^{\ast}).
\end{equation}
The online network selects; the target network evaluates.
This reduces the noise amplification from using the same noisy Q for both roles.

\textbf{Dueling DQN} decomposes the Q-function structurally:
\begin{equation}
Q(s,a) = V(s) + A(s,a),
\end{equation}
where $V(s)$ captures overall state value and $A(s,a)$ captures the relative advantage of each action.
By isolating the state-value component, estimation becomes more stable.

\begin{keyidea}[DQN's Legacy]
The DQN family retains Qmax's aggressive optimization but puts reins on it.
The Bellman equation is still the primary objective---but it now requires heavy engineering stabilization.
\end{keyidea}

% ===========================================================
\section{Stage 6: Policy Gradient---The Philosophical Watershed}
\label{sec:stage6-pg}
% ===========================================================

Q-learning is fundamentally \textbf{value-centric}: it posits an objectively optimal $Q^{\ast}(s,a)$ and approaches it through Bellman consistency.
Policy Gradient (Williams 1992; Sutton et al.\ 1999) is entirely different.
It directly optimizes the policy $\pi_\theta(a|s)$.

The objective:
\begin{equation}
J(\theta) = \mathbb{E}_{\tau \sim \pi_\theta}\bigl[R(\tau)\bigr].
\end{equation}
The gradient (via the likelihood ratio trick):
\begin{equation}
\label{eq:pg}
\nabla_\theta J(\theta) = \mathbb{E}\Bigl[\nabla_\theta \log \pi_\theta(a|s)\, A(s,a)\Bigr].
\end{equation}

The meaning is intuitive: if an action's advantage is positive, increase its probability; if negative, decrease it.
This is not hard $\max$---it is \textbf{probability reshaping}.

\subsection{Two Layers of Probability}

The statistical essence of policy gradient can be understood through two layers:
\begin{enumerate}
    \item \textbf{Rollout probability}: which actions get sampled depends on $\pi_\theta(a|s)$.
    \item \textbf{Update weighting}: sampled actions are re-weighted by advantage $A(s,a)$.
\end{enumerate}
Policy gradient does not simply copy sampled behaviors.
Within the actual sampling distribution, it re-weights behaviors by reward as preferences---closely mirroring statistical estimation.

\subsection{A Change of Objective}

The shift from Q-learning to policy gradient is not merely algorithmic---the optimization objective changes.

Q-learning's Bellman loss is a \textbf{consistency loss}:
\begin{equation}
\mathcal{L}_Q = \bigl(Q(s,a) - r - \gamma\,\max_{a'} Q(s',a')\bigr)^2.
\end{equation}
It pursues value-function self-consistency---fitting ``true'' values.

Policy gradient's loss is \textbf{behavior distribution shaping}:
\begin{equation}
\mathcal{L}_\pi = -\log \pi(a|s)\, A(s,a).
\end{equation}
It pursues increasing good-behavior probabilities and decreasing bad ones.

\begin{keyidea}[The Watershed]
Q-learning says: ``The action with the highest estimated value should be selected.''
Policy gradient says: ``Behaviors that perform statistically better should occur more frequently.''
This is the turn from god's-eye value solving to sampling-based behavioral preference shaping.
\end{keyidea}

\subsection{Policy Gradient as Weighted Maximum Likelihood}

The connection to statistics is precise.
Ordinary maximum likelihood (behavior cloning) maximizes $\sum_i \log \pi_\theta(a_i|s_i)$---every observed action is reinforced equally.
Policy gradient maximizes a \emph{weighted} log-likelihood:
\begin{equation}
\sum_i A(s_i, a_i)\,\log \pi_\theta(a_i|s_i).
\end{equation}
MLE models behavioral frequency; policy gradient models behavioral frequency with preference weights.

% ===========================================================
\section{Stage 7: Actor-Critic and the PPO/SAC Era (2015--2018)}
\label{sec:stage7-ac}
% ===========================================================

Policy gradient needs advantage estimates.
Actor-Critic methods maintain both a policy (actor) and a value function (critic):

\begin{table}[ht]
\centering
\caption{Critic variants in Actor-Critic architectures.}
\label{tab:critic-types}
\begin{tabular}{lll}
\toprule
\textbf{Critic Type} & \textbf{What It Learns} & \textbf{Typical Algorithms} \\
\midrule
V-critic    & $V(s)$    & A2C, A3C, PPO \\
Q-critic    & $Q(s,a)$  & DDPG, TD3, SAC \\
\bottomrule
\end{tabular}
\end{table}

PPO-style methods use a V-critic to estimate advantage via TD residuals or GAE:
\begin{equation}
\hat{A}_t \approx r_t + \gamma\, V(s_{t+1}) - V(s_t).
\end{equation}
The critic here is a statistical control variate---not an attempt to find a globally optimal $Q^{\ast}$.

\subsection{The Key Turn: Policy Becomes Primary}

TRPO, PPO, and SAC marked a fundamental shift.
The optimization objective is no longer Bellman consistency but:
\begin{equation}
\max_\theta\; \mathbb{E}_{\pi_\theta}\bigl[R(\tau)\bigr].
\end{equation}
The Bellman equation is \emph{demoted to an auxiliary role}---the critic uses it to estimate advantage, but it is no longer the primary optimization formula.

\subsection{Greedy with Brakes}

Modern methods can be understood as \textbf{greedy optimization with brakes}:

\begin{itemize}
    \item \textbf{Why greedy?}  Data is scarce.  Without greedy amplification, high-reward signals propagate too slowly.
    \item \textbf{Why brakes?}  Greedy amplifies noise, bias, and function approximation error.
\end{itemize}

\begin{table}[ht]
\centering
\caption{How modern algorithms balance greediness and stability.}
\label{tab:brakes}
\begin{tabular}{ll}
\toprule
\textbf{Algorithm / Technique} & \textbf{Essential Function} \\
\midrule
PPO clipping              & Limit overly aggressive policy updates \\
TRPO / KL constraint      & Trust region prevents policy jumps \\
Entropy regularization    & Prevent premature deterministic collapse \\
TD3 clipped double Q      & Conservative estimates suppress overestimation \\
SAC                       & Soft greedy + entropy + Q-critic \\
Target network            & Reduce bootstrap target drift \\
Double Q                  & Reduce $\max$ overestimation \\
\bottomrule
\end{tabular}
\end{table}

The shared principle: \emph{do not abandon optimization, but prevent the optimizer from exploiting estimation noise.}

\subsection{V-Only as a Relaxation of Qmax}

PPO replaces hard $\arg\max$ with soft advantage-weighted updates:
\begin{align}
\text{Q-learning:} \quad & a^{\ast} = \arg\max_a Q(s,a), \\
\text{PPO:} \quad & \nabla_\theta \log \pi_\theta(a|s)\, A(s,a).
\end{align}
The former is hard optimization; the latter is soft preference adjustment.
PPO is more stable but not ``more globally optimal.''
What it pursues is: stably pushing the policy toward better directions under limited sampling.

% ===========================================================
\section{Stage 8: RLHF---Preferences Replace Engineered Rewards (2017--2022)}
\label{sec:stage8-rlhf}
% ===========================================================

Christiano et al.\ (2017) proposed learning a reward model from human preference comparisons.
Ouyang et al.\ (2022, InstructGPT) applied it at scale to LLMs.

The shift: reward is no longer a hand-designed function but is \emph{sampled from human preferences}---noisy, contextual, local, sometimes contradictory.
Forcing a globally consistent Bellman-optimal $Q^{\ast}$ onto such preferences may itself be a fiction.

However, RLHF's policy training still used PPO internally.
PPO's critic still relied on the Bellman equation for value estimation.
Preferences changed the \emph{source} of reward, but the Bellman equation had not yet fully exited the optimization mechanism.

% ===========================================================
\section{Stage 9: DPO and GRPO---The Bellman Equation Disappears (2023--2025)}
\label{sec:stage9-dpo}
% ===========================================================

This is the most dramatic rupture in the entire arc.

\textbf{DPO} (Rafailov et al., 2023) directly optimizes the policy on preference pairs, eliminating the reward model and the RL training loop entirely.
No critic, no value network, no Bellman bootstrap.

\textbf{GRPO} (DeepSeek, 2025) goes further: for each prompt, generate a group of responses and use within-group relative ranking as the baseline to replace the critic.
It is pure policy gradient plus statistical preference re-weighting.

\begin{keyidea}[The Complete Break]
In DPO/GRPO there is:
\begin{itemize}
    \item No value network ($V$ or $Q$).
    \item No Bellman consistency loss.
    \item No bootstrap target.
\end{itemize}
The optimization objective is pure preference distribution shaping.
The Bellman equation is completely gone.
\end{keyidea}

This is not regression---it is the inevitable endpoint of a trend.
In a world where human preferences are inherently noisy, contextual, and ad-hoc, maintaining a globally consistent Bellman fixed point is neither necessary nor natural.
Statistical preference optimization is the correct response to this reality.

% ===========================================================
\section{The Diminishing Role of the Bellman Equation}
\label{sec:bellman-role}
% ===========================================================

The entire nine-stage evolution can be summarized by tracking a single variable: \emph{how central is the Bellman equation to the optimization?}

\begin{table}[ht]
\centering
\caption{The Bellman equation's role at each stage.}
\label{tab:bellman-role}
\begin{tabular}{lll}
\toprule
\textbf{Stage} & \textbf{Era} & \textbf{Bellman's Role} \\
\midrule
1--4. DP / Q-learning       & 1957--1989 & The entire optimization \\
5. DQN family               & 2013--2016 & The entire optimization (+ stability tricks) \\
6. Early Actor-Critic       & $\sim$2000 & Primary objective (policy is a means) \\
7. PPO / SAC                & 2015--2018 & Auxiliary tool (critic uses it for advantage) \\
8. RLHF                     & 2017--2022 & Residual (inside PPO's critic) \\
9. DPO / GRPO               & 2023--2025 & Completely absent \\
\bottomrule
\end{tabular}
\end{table}

% ===========================================================
\section{The $\arg\max$ Distinction: Conclusion vs.\ Process}
\label{sec:argmax-distinction}
% ===========================================================

A subtle but critical point underlies the entire arc.

In Bellman optimality, $\arg\max$ is a \textbf{conclusion}:
\emph{if} you already know the true $Q^{\ast}$, \emph{then} choosing $\arg\max Q^{\ast}$ is correct.

During training, $\arg\max$ becomes the \textbf{process} at every step:
even though $Q$ is noisy, you forcefully select the current maximum.

These are completely different operations.
An analogy: Bellman optimality's $\arg\max$ is knowing that daily exercise is good \emph{when you are already fit}.
Hard greedy during training is attempting maximum-intensity exercise every day \emph{when your body is in poor condition}---overshooting, injury, collapse.

Modern RL makes the optimization ``softer'':
\begin{itemize}
    \item PPO uses clipping to prevent the policy from changing too much in one step.
    \item SAC uses entropy to prevent premature collapse.
    \item Double Q / TD3 uses conservative estimates to suppress overestimation.
    \item RLHF uses KL regularization to stay close to the reference policy.
\end{itemize}

None of these deny optimality.
They acknowledge that under limited data, you cannot treat theoretical optimality conditions as hard execution rules at every training step.

% ===========================================================
\section{From God's-Eye View to Statistical View}
\label{sec:god-to-stat}
% ===========================================================

Classical DP adopts a god's-eye view: there exists a true model, true transition probabilities, a true reward function, and true optimal $Q^{\ast}$, $V^{\ast}$.
The algorithm's task is to compute these true values.

The statistical view is different.
There is no directly observable truth---only samples, assumptions, likelihoods, and estimators.
Maximum likelihood estimation does not seek cosmic truth; it asks:
\emph{what parameters best explain the current observed data?}

As RL has evolved, it has increasingly adopted this statistical worldview:
\begin{itemize}
    \item Policy gradient, RLHF, and GRPO do not insist on a globally consistent $Q^{\ast}$.
    \item They care about: which behaviors brought better rewards in sampling statistics, and therefore should occur more frequently.
\end{itemize}

This is the shift from ``finding true values'' to ``shaping behavioral distributions.''

\begin{keyidea}[Two Worldviews]
DP is exact optimization under a known model---\textbf{analytical optimal control}.\\
Modern RL is recursive preference optimization under noise, bias, and distribution drift---\textbf{statistical preference learning}.
\end{keyidea}

% ===========================================================
\section{The Six Phases at a Glance}
\label{sec:six-phases}
% ===========================================================

The nine stages cluster naturally into six conceptual phases.

\subsection*{Phase 1: DP / V-Function}
\emph{Goal}: stable evaluation and globally optimal solving under known models.\\
\emph{Character}: exact expectations, clean Bellman fixed point, stable V-function.\\
\emph{Limitation}: action optimization is not direct.

\subsection*{Phase 2: Q-Learning / Qmax}
\emph{Goal}: explicit action optimization without a model.\\
\emph{Character}: high data efficiency, off-policy, fast greedy propagation.\\
\emph{Limitation}: severe overestimation, instability.

\subsection*{Phase 3: DQN Family}
\emph{Goal}: scale Q-learning with deep networks.\\
\emph{Character}: still value-centric, still Bellman-centered, but with heavy engineering stabilization.\\
\emph{Limitation}: $\max$ bias persists despite repairs.

\subsection*{Phase 4: Policy Gradient / Actor-Critic}
\emph{Goal}: directly optimize the policy distribution.\\
\emph{Character}: shift from value-centric to behavior-centric; no longer insisting on perfect Q.\\
\emph{Limitation}: high variance, lower data efficiency.

\subsection*{Phase 5: PPO / SAC / TD3}
\emph{Goal}: fuse value and policy, balance greediness and stability.\\
\emph{Character}: constrained/regularized policy updates; the Bellman equation becomes auxiliary.\\
\emph{Limitation}: still requires careful hyperparameter tuning.

\subsection*{Phase 6: RLHF / DPO / GRPO}
\emph{Goal}: align with human preferences.\\
\emph{Character}: reward from preferences (not engineering); algorithms resemble statistical preference learning.\\
\emph{Essence}: the goal is not absolute optimality, but placing the behavioral distribution within the zone of human preferences.

% ===========================================================
\section{Chapter Summary}
\label{sec:ch2-summary}
% ===========================================================

Ten points compress the grand arc:

\begin{enumerate}
    \item DP is exact optimization under known models; RL is sampling-based optimization under unknown models.
    \item The Bellman equation describes an ideal fixed point; the training process is noisy nonlinear recursive iteration.
    \item V is stable but conservative; Q enables explicit action optimization but more easily amplifies errors.
    \item $\max$ is data-efficient because it rapidly propagates high values---but it also preferentially selects positive noise.
    \item Modern RL's core is not pursuing global optimality, but learning stable, useful, statistically superior policies under limited data.
    \item Policy Gradient is the watershed: from finding the optimal value function to directly shaping behavior distributions.
    \item Actor-Critic fuses value estimation and policy optimization---not simply ``using Q'' or ``not using Q.''
    \item PPO / SAC / TD3 are greedy optimization with brakes.
    \item RLHF / GRPO demonstrate that real-world preferences are inherently ad-hoc and statistical---not suited for a globally Bellman-consistent value function.
    \item Modern RL more closely resembles statistical preference learning than classical analytical optimal control.
\end{enumerate}

The entire thread in one powerful statement:

\begin{keyidea}[The Essence]
DP pursues the global optimal solution.
RL pursues evolving good-enough behavioral preferences from limited, noisy, biased data.
The field's history oscillates between ``too much averaging---can't learn'' and ``too much greediness---collapses.''
The answer is not choosing one side, but:
use greediness to amplify effective signals; use statistical constraints and regularization to prevent the system from being carried away by noise.
\end{keyidea}


\part{The Mechanisms}
\chapter{The Loss Landscape of RL}
\label{ch:loss-landscape}

Reinforcement learning borrows the optimization machinery of supervised learning---forward pass, loss computation, backpropagation, parameter update---but the losses it optimizes are fundamentally different. In supervised learning, the loss measures deviation from externally provided labels on a fixed dataset. In RL, the ``labels'' are constructed by the agent itself, the data distribution shifts as the policy changes, and the loss must encode not just fitting accuracy but also exploration incentives and distributional control.

This chapter traces the evolution of RL loss functions: from MSE regression on Bellman targets, through direct Q-maximization, to weighted log-likelihood and entropy regularization. Each transition addresses a specific failure mode of the previous formulation.

%% ----------------------------------------------------------------
\section{Value-Based Loss: MSE on Self-Generated Labels}
\label{sec:td-mse}

The starting point is deceptively simple. DQN and Q-learning frame RL as a regression problem:
\begin{equation}
  \mathcal{L}_{\text{critic}} = \bigl(Q_\theta(s,a) - y\bigr)^2,
  \qquad
  y = r + \gamma \max_{a'} Q_{\bar\theta}(s', a'),
  \label{eq:td-loss}
\end{equation}
where $\bar\theta$ denotes the parameters of a slowly updated target network.

\begin{keyidea}[Self-Bootstrapping Regression]
The target $y$ is not an external ground-truth label. It is computed by the model itself through the Bellman bootstrap and the $\max$ operator. The model creates its own labels, then fits those labels. This is the core characteristic of TD learning and the reason it is \emph{not} ordinary MSE.
\end{keyidea}

Two consequences follow immediately:
\begin{enumerate}
  \item \textbf{The target is noisy.} Because $Q_{\bar\theta}$ is an approximation, $y$ carries estimation error that feeds back into training.
  \item \textbf{The $\max$ operator amplifies noise.} Among noisy estimates, the maximum is biased upward (Section~\ref{sec:overestimation}).
\end{enumerate}

%% ----------------------------------------------------------------
\section{Q Overestimation}
\label{sec:overestimation}

Suppose each action's Q-estimate decomposes as
\begin{equation}
  \hat{Q}(a) = Q^*(a) + \epsilon(a),
  \qquad \epsilon(a) \sim \text{noise}.
\end{equation}
When we take the maximum over actions, the action whose noise $\epsilon(a)$ happens to be largest is preferentially selected. By Jensen's inequality applied to the convex $\max$ function:
\begin{equation}
  \mathbb{E}\!\bigl[\max_a \hat{Q}(a)\bigr]
  \;\ge\;
  \max_a \mathbb{E}\!\bigl[\hat{Q}(a)\bigr].
  \label{eq:overestimation}
\end{equation}

This is not a peculiarity of DQN. \emph{Any} algorithm that relies on $\max Q$, $\arg\max Q$, or advantage ranking faces the same issue. The difference between algorithms lies in how aggressively they mitigate it.

\begin{keyidea}[Overestimation Is Structural]
Q overestimation is not a bug in a particular implementation. It is a mathematical consequence of taking the maximum of noisy estimates. The more actions, the larger the bias.
\end{keyidea}

%% ----------------------------------------------------------------
\section{Double DQN and TD3: Stabilizing the Target}
\label{sec:double-td3}

\subsection{Double DQN}

The core idea is to \emph{decouple action selection from value evaluation}. Standard DQN uses the same (noisy) network for both:
\begin{equation}
  y_{\text{DQN}} = r + \gamma \max_{a'} Q_{\bar\theta}(s', a').
\end{equation}
Double DQN selects the action with the online network but evaluates it with the target network:
\begin{equation}
  a^* = \arg\max_{a'} Q_\theta(s', a'),
  \qquad
  y_{\text{DDQN}} = r + \gamma\, Q_{\bar\theta}(s', a^*).
  \label{eq:double-dqn}
\end{equation}
This does not eliminate bias entirely, but it breaks the positive feedback loop in which the same noisy Q is responsible for both choosing and scoring.

\subsection{TD3: Clipped Double-Q for Continuous Actions}

TD3 extends the idea to continuous action spaces. It maintains two independent critics $Q_{\phi_1}$, $Q_{\phi_2}$ and uses the more pessimistic estimate:
\begin{equation}
  y = r + \gamma \min\bigl(Q_{\phi_1}(s', \tilde{a}'),\; Q_{\phi_2}(s', \tilde{a}')\bigr),
  \qquad
  \tilde{a}' = \mu_{\bar\theta}(s') + \text{clip}(\epsilon, -c, c).
  \label{eq:td3-target}
\end{equation}
The $\min$ produces a conservative estimate that counteracts the actor's tendency to exploit critic errors. Two additional mechanisms---\emph{delayed policy updates} (let the critic stabilize before updating the actor) and \emph{target policy smoothing} (add clipped noise to the target action)---further prevent the actor from chasing sharp, unreliable peaks in the Q landscape.

%% ----------------------------------------------------------------
\section{Max-Q Loss: The Actor Directly Maximizes Q}
\label{sec:max-q}

In discrete action spaces, $\arg\max_a Q(s,a)$ is trivial: enumerate all actions, pick the winner. In continuous spaces this is intractable, so DDPG introduces an actor network $\mu_\theta(s)$ trained to output the action that maximizes the critic:
\begin{equation}
  \mathcal{L}_{\text{actor}} = -\,Q_\phi\!\bigl(s,\, \mu_\theta(s)\bigr).
  \label{eq:ddpg-actor}
\end{equation}
Minimizing this loss is equivalent to gradient ascent on $Q$ with respect to the actor's output. The gradient signal flows through the critic: $\nabla_\theta \mathcal{L}_{\text{actor}} = -\nabla_a Q_\phi(s,a)\big|_{a=\mu_\theta(s)} \cdot \nabla_\theta \mu_\theta(s)$.

\begin{keyidea}[Differentiable Argmax]
DDPG's actor loss turns the discrete, non-differentiable $\arg\max$ operation into a differentiable gradient optimization problem. The actor uses the critic's gradient landscape $\nabla_a Q$ to steer its output toward high-value actions.
\end{keyidea}

\subsection{The Fundamental Risk}

The critic's Q is not the true value function---it has noise, bias, and overestimation. If the actor relentlessly climbs the Q surface, it may converge not to truly good actions but to actions where the critic's \emph{estimation error} is largest. This is why TD3 requires its suite of protective measures (double critic, delayed updates, target smoothing).

\subsection{Placing Max-Q in the Loss Evolution}

\begin{center}
\begin{tabular}{@{} l l @{}}
\toprule
\textbf{Loss type} & \textbf{What it does} \\
\midrule
MSE (TD) & Make $Q$ approximate a Bellman target --- passive fitting \\
Max-Q (DDPG actor) & Make the actor output high-$Q$ actions --- actively exploiting the Q landscape \\
Policy gradient & Make good actions more probable --- directly shaping the distribution \\
\bottomrule
\end{tabular}
\end{center}

MSE says ``I estimate value.'' Max-Q says ``I exploit my estimates for decisions.'' Policy gradient says ``I directly optimize decisions.''

%% ----------------------------------------------------------------
\section{Policy Gradient: Gradient First, Loss Second}
\label{sec:pg-loss}

Policy gradient methods bypass the Q landscape entirely. The objective is the expected return:
\begin{equation}
  J(\theta) = \mathbb{E}_{\pi_\theta}\!\bigl[\text{return}\bigr].
\end{equation}
The policy gradient theorem yields:
\begin{equation}
  \nabla_\theta J(\theta)
  = \mathbb{E}_{s,a\sim\pi_\theta}\!\bigl[
    \nabla_\theta \log\pi_\theta(a\mid s)\;\hat{A}(s,a)
  \bigr],
  \label{eq:pg}
\end{equation}
where $\hat{A}(s,a)$ is an advantage estimate.

\subsection{Why ``Gradient First''}

A common source of confusion: value-based methods start from a loss and derive a gradient, but policy gradient derivations start from the gradient and only mention a loss at the end. The reason is structural.

\textbf{Value-based:} There is a natural regression target $y$. The loss $\mathcal{L} = (Q - y)^2$ is the starting point; the gradient follows from differentiation.

\textbf{Policy gradient:} The starting point is the objective $J(\theta)$. The gradient~\eqref{eq:pg} is derived via the likelihood ratio trick applied to trajectory probabilities. No loss function is written first.

To use automatic differentiation frameworks, we \emph{construct a surrogate loss after the fact}:
\begin{equation}
  \mathcal{L}_{\text{PG}} = -\log\pi_\theta(a\mid s)\;\hat{A}(s,a).
  \label{eq:pg-loss}
\end{equation}
Differentiating this with respect to $\theta$ recovers exactly the policy gradient~\eqref{eq:pg} (with the sign flip because optimizers minimize). The loss is an engineering artifact for autograd, not the conceptual starting point.

\begin{keyidea}[Constructed Loss]
In value-based RL, loss comes first and gradient follows. In policy gradient RL, the gradient comes first from mathematical derivation; the loss is constructed afterward so that \texttt{loss.backward()} produces the correct gradient.
\end{keyidea}

Yet in code, both families share the same training loop: \texttt{loss = compute\_loss(batch); loss.backward(); optimizer.step()}.

%% ----------------------------------------------------------------
\section{Connection to Maximum Likelihood}
\label{sec:mle-connection}

\subsection{Behavior Cloning as Unweighted MLE}

Given demonstration data $\{(s_i, a_i)\}$, behavior cloning maximizes:
\begin{equation}
  \sum_i \log\pi_\theta(a_i \mid s_i).
\end{equation}
Every observed action receives equal weight. The model learns ``a human did this,'' with no notion of whether ``this'' was good or bad.

\subsection{Policy Gradient as Weighted MLE}

Policy gradient adds a quality signal:
\begin{equation}
  \mathcal{L}_{\text{PG}} = -\hat{A}(s,a)\;\log\pi_\theta(a \mid s).
\end{equation}
The advantage $\hat{A}$ serves as a per-sample weight. Actions that led to higher returns get amplified; actions that underperformed get suppressed.

\begin{keyidea}[From MLE to Weighted MLE]
Behavior cloning (unweighted MLE) models behavioral frequency: ``this action appeared, so increase its probability.'' Policy gradient (weighted MLE) models behavioral frequency with preference weights: ``this action appeared \emph{and} brought high returns, so increase its probability proportionally.''
\end{keyidea}

This perspective makes the conceptual shift transparent:
\begin{center}
\begin{tabular}{@{} l l @{}}
\toprule
\textbf{Method} & \textbf{What it models} \\
\midrule
MLE / Behavior Cloning & Frequency of observed actions \\
Policy Gradient & Frequency weighted by quality of outcomes \\
\bottomrule
\end{tabular}
\end{center}

%% ----------------------------------------------------------------
\section{Entropy: Sampling Distribution Regularization}
\label{sec:entropy}

PPO and SAC both include an entropy bonus $\mathcal{H}[\pi_\theta(\cdot \mid s)]$. Its role is often trivialized as ``encouraging randomness.'' The real function is more precise: it prevents the policy distribution from collapsing prematurely.

\subsection{The Collapse Mechanism}

Without entropy regularization, the policy gradient pushes probability mass toward actions with high estimated advantage. The resulting feedback loop is:
\begin{equation*}
  \text{policy concentrates}
  \;\to\; \text{sampling narrows}
  \;\to\; \text{data diversity drops}
  \;\to\; \text{Q/advantage estimates degrade}
  \;\to\; \text{local optima.}
\end{equation*}

\subsection{Entropy as a Distributional Force}

Policy optimization involves two competing forces:

\begin{center}
\begin{tabular}{@{} l l @{}}
\toprule
\textbf{Force} & \textbf{Effect} \\
\midrule
Preference term $\hat{A}\cdot\log\pi$ & Push probability toward high-advantage actions \\
Entropy term $\mathcal{H}[\pi]$ & Pull the distribution toward uniformity \\
\bottomrule
\end{tabular}
\end{center}

The learned policy is the equilibrium between these two forces---neither a uniform distribution nor a collapsed delta, but a controlled trade-off.

\subsection{SAC: Entropy as Part of the Objective}

SAC goes further by writing the objective as:
\begin{equation}
  J_{\text{SAC}}(\theta) = \mathbb{E}\!\Bigl[\sum_t r_t + \alpha\,\mathcal{H}\bigl[\pi_\theta(\cdot \mid s_t)\bigr]\Bigr],
  \label{eq:sac-obj}
\end{equation}
where $\alpha$ is a temperature parameter (often auto-tuned). Exploration is no longer an external trick appended to the loss; it is a first-class component of the optimization objective. Sampling diversity is something the agent is \emph{rewarded} for.

%% ----------------------------------------------------------------
\section{The Full Loss Evolution}
\label{sec:loss-evolution}

We can now see the complete arc:

\begin{center}
\begin{tabular}{@{} l l l @{}}
\toprule
\textbf{Loss layer} & \textbf{Formula sketch} & \textbf{What it controls} \\
\midrule
MSE (TD) & $(Q - y)^2$ & Value estimation accuracy \\
Max-Q (actor) & $-Q(s, \mu_\theta(s))$ & Action quality via critic landscape \\
Weighted $\log\pi$ & $-\hat{A}\cdot\log\pi_\theta(a \mid s)$ & Preference-weighted behavior \\
Entropy & $+\alpha\,\mathcal{H}[\pi]$ & Sampling diversity control \\
\bottomrule
\end{tabular}
\end{center}

These four layers correspond to four progressively deeper concerns:
\begin{enumerate}
  \item \textbf{Estimate value} (how much is this worth?).
  \item \textbf{Exploit estimates} (given the value landscape, what should I do?).
  \item \textbf{Shape the policy} (make good actions more probable).
  \item \textbf{Control the distribution} (keep the policy from collapsing).
\end{enumerate}

Modern algorithms like SAC combine all four layers---a critic trained with TD-MSE, an actor that maximizes entropy-augmented Q, and an entropy term that regularizes the policy distribution---into a single coherent framework. Understanding each layer in isolation makes the composite loss transparent.

\chapter{On-Policy, Off-Policy, and the Sampling Problem}
\label{ch:on-off-policy}

Whether an algorithm is on-policy or off-policy is not a label to memorize per algorithm name. It is a consequence of a single question about the algorithm's update rule. This chapter derives the answer from first principles for every major algorithm family, then examines the practical implications for data efficiency.

%% ----------------------------------------------------------------
\section{The First-Principles Criterion}
\label{sec:criterion}

To classify any RL algorithm, ask one question:

\begin{keyidea}[The Sampling Criterion]
Does the update formula require the actions in the training data to have been sampled by the \emph{current} policy $\pi_\theta$?
\end{keyidea}

More precisely, check whether two chains are consistent:
\begin{itemize}
  \item \textbf{Sampling chain:} Which policy selected the actions in the data?
  \item \textbf{Optimization chain:} Which policy's parameters are being updated?
\end{itemize}

This yields three categories:
\begin{center}
\begin{tabular}{@{} l l @{}}
\toprule
\textbf{Category} & \textbf{Requirement on data actions} \\
\midrule
On-policy & $a \sim \pi_\theta(a \mid s)$ \quad (current policy) \\
Off-policy & $a \sim \mu(a \mid s)$ \quad (any behavior policy) \\
Near on-policy & $a \sim \pi_{\text{old}}(a \mid s)$, with $\pi_{\text{old}} \approx \pi_\theta$ enforced \\
\bottomrule
\end{tabular}
\end{center}

The real criterion is not ``whether there is a replay buffer'' but whether the action distribution in the data must match the distribution of the policy currently being optimized.

%% ----------------------------------------------------------------
\section{Why Policy Gradient Is On-Policy}
\label{sec:pg-on}

The policy gradient takes the form:
\begin{equation}
  \nabla_\theta J(\theta)
  = \mathbb{E}_{s,a \,\sim\, \pi_\theta}\!\bigl[
      \nabla_\theta \log\pi_\theta(a \mid s)\;\hat{A}^{\pi_\theta}(s,a)
    \bigr].
  \label{eq:pg-onpolicy}
\end{equation}

The expectation is taken over $s, a \sim \pi_\theta$. This is not a notational convenience; it is a mathematical requirement. The derivation uses the identity $\nabla_\theta \pi_\theta(a \mid s) = \pi_\theta(a \mid s)\,\nabla_\theta\log\pi_\theta(a \mid s)$, which only produces an unbiased gradient estimate when actions are drawn from $\pi_\theta$.

If actions came from an old or different policy $\mu \neq \pi_\theta$, the weighted log-likelihood $\log\pi_\theta(a \mid s)\cdot\hat{A}(s,a)$ is no longer an unbiased estimate of $\nabla_\theta J(\theta)$.

\begin{keyidea}[PG Is On-Policy by Derivation]
REINFORCE, A2C, and A3C are on-policy not because of their names, but because the gradient formula~\eqref{eq:pg-onpolicy} requires $a$ to come from the current $\pi_\theta$'s sampling distribution.
\end{keyidea}

%% ----------------------------------------------------------------
\section{Why Q-Learning Is Off-Policy}
\label{sec:ql-off}

The Q-learning target is:
\begin{equation}
  y = r + \gamma \max_{a'} Q_{\bar\theta}(s', a').
  \label{eq:ql-target}
\end{equation}

Crucially, this target does not depend on how action $a$ was selected. The transition $(s, a, r, s')$ tells us: ``someone executed $a$ in $s$ and observed $r, s'$.'' The target then asks: ``if the agent acts greedily from $s'$ onward, what is the expected return?'' The identity of ``someone'' is irrelevant---it could be the current policy, an old policy, a human operator, or a random controller.

\begin{keyidea}[Bellman Doesn't Care Who Sampled]
Q-learning's target involves $\max_{a'} Q(s', a')$, which is independent of the behavior policy that generated the data. This is why Q-learning can learn from a replay buffer of transitions collected by arbitrary policies---the fundamental reason it is off-policy.
\end{keyidea}

The same logic extends to DQN and Double DQN: the Bellman target depends on the \emph{target network's} value estimates, not on the policy that produced the stored transitions.

%% ----------------------------------------------------------------
\section{Why PPO Is Near On-Policy}
\label{sec:ppo-near}

PPO collects data with a recent snapshot of the policy $\pi_{\text{old}}$, then performs multiple gradient steps on the current policy $\pi_\theta$. This is not purely on-policy (the data was not sampled by $\pi_\theta$) nor purely off-policy (PPO cannot use arbitrary historical data).

PPO corrects for the distribution mismatch using an importance ratio:
\begin{equation}
  r_t(\theta) = \frac{\pi_\theta(a_t \mid s_t)}{\pi_{\text{old}}(a_t \mid s_t)},
\end{equation}
and then clips it to prevent large updates:
\begin{equation}
  \mathcal{L}_{\text{PPO}} = -\min\!\bigl(
    r_t(\theta)\,\hat{A}_t,\;
    \operatorname{clip}(r_t(\theta),\, 1{-}\epsilon,\, 1{+}\epsilon)\,\hat{A}_t
  \bigr).
  \label{eq:ppo-clip}
\end{equation}

The clipping enforces that $\pi_\theta$ and $\pi_{\text{old}}$ cannot diverge too much. Data from the most recent old policy is acceptable; data from a much older policy is not.

\begin{keyidea}[Near On-Policy = Bounded Staleness]
PPO allows limited reuse of recent data, but constrains distribution shift through ratio clipping. It is neither on-policy (data is slightly stale) nor off-policy (it cannot use arbitrary historical data). The precise term is \emph{near on-policy}.
\end{keyidea}

TRPO follows the same logic but enforces the constraint via a KL-divergence trust region rather than clipping.

%% ----------------------------------------------------------------
\section{Why DDPG, TD3, and SAC Are Off-Policy}
\label{sec:offpolicy-ac}

All three algorithms use a replay buffer, but the deeper reason for their off-policy nature is in their update rules.

\subsection{Critic Update}

The critic targets have the same structure as Q-learning---a Bellman backup that does not depend on the behavior policy:
\begin{align}
  y_{\text{DDPG}} &= r + \gamma\, Q_{\bar\phi}\!\bigl(s',\, \mu_{\bar\theta}(s')\bigr), \\
  y_{\text{TD3}}  &= r + \gamma\, \min_{i=1,2} Q_{\bar\phi_i}\!\bigl(s',\, \mu_{\bar\theta}(s') + \tilde\epsilon\bigr), \\
  y_{\text{SAC}}  &= r + \gamma\, \min_{i=1,2} Q_{\bar\phi_i}(s', \tilde{a}') - \alpha\log\pi_\theta(\tilde{a}' \mid s'),
  \quad \tilde{a}' \sim \pi_\theta(\cdot \mid s').
\end{align}

In each case, the action $a$ stored in the replay buffer can come from any previous version of the policy. The target is constructed from target networks and does not require $a \sim \pi_\theta$.

\subsection{Actor Update}

The actor optimizes \emph{through} the critic:
\begin{equation}
  \nabla_\theta J \approx \nabla_\theta\, Q_\phi\!\bigl(s,\, \mu_\theta(s)\bigr)
  \quad\text{(DDPG/TD3)},
  \qquad
  \nabla_\theta\, \mathbb{E}_{a\sim\pi_\theta}\!\bigl[Q_\phi(s,a) - \alpha\log\pi_\theta(a \mid s)\bigr]
  \quad\text{(SAC)}.
\end{equation}
The states $s$ are drawn from the replay buffer, but the actions used in the gradient are freshly computed by the current actor/policy. The actor does not need the stored action to match $\pi_\theta$---it generates its own actions at update time.

%% ----------------------------------------------------------------
\section{Quick Reference Table}
\label{sec:reference-table}

\begin{center}
\begin{tabular}{@{} l l l l @{}}
\toprule
\textbf{Algorithm} & \textbf{Data source} & \textbf{Requires $a \sim \pi_\theta$?} & \textbf{Verdict} \\
\midrule
REINFORCE   & Current policy       & Yes          & On-policy \\
A2C / A3C   & Current policy       & Yes          & On-policy \\
PPO         & Recent old policy    & Approximately & Near on-policy \\
TRPO        & Recent old policy    & Approximately & Near on-policy \\
Q-learning  & Any behavior policy  & No           & Off-policy \\
DQN         & Replay buffer        & No           & Off-policy \\
Double DQN  & Replay buffer        & No           & Off-policy \\
DDPG        & Replay buffer        & No           & Off-policy \\
TD3         & Replay buffer        & No           & Off-policy \\
SAC         & Replay buffer        & No           & Off-policy \\
\bottomrule
\end{tabular}
\end{center}

%% ----------------------------------------------------------------
\section{Data Efficiency: Theory vs.\ Engineering}
\label{sec:data-efficiency}

\subsection{Theoretical Sample Efficiency}

Off-policy methods can replay each transition many times, giving them a clear theoretical advantage in sample efficiency. A single $(s, a, r, s')$ tuple collected months ago can still contribute to today's gradient update.

On-policy methods discard data after each policy update. Once $\pi_\theta$ changes, old trajectories no longer provide unbiased gradient estimates, so the data must be thrown away.

\subsection{Engineering Efficiency Is a Different Story}

In practice, raw sample efficiency is only one factor:

\begin{center}
\begin{tabular}{@{} l p{5.5cm} p{5.5cm} @{}}
\toprule
\textbf{Factor} & \textbf{PPO (near on-policy)} & \textbf{SAC (off-policy)} \\
\midrule
Sample efficiency & Lower---data discarded after a few epochs & Higher---replay buffer reused extensively \\
Training stability & Very stable; robust to hyperparameters & Can diverge if Q-estimates blow up \\
Hyperparameter tuning & Easy; few sensitive knobs & Harder; temperature $\alpha$, replay ratio, critic architecture matter \\
Parallelism & Excellent; scales to thousands of workers & Moderate; replay buffer is a bottleneck \\
Reproducibility & High & Lower; sensitive to seeds and ordering \\
\bottomrule
\end{tabular}
\end{center}

\begin{keyidea}[Theoretical $\neq$ Engineering Efficiency]
Theoretical sample efficiency measures how many times each transition can be reused by the optimizer. Engineering efficiency measures whether, given the same engineering effort, you can stably obtain a good policy. PPO wastes samples but is easy to make work. SAC reuses samples but demands more careful tuning. If the off-policy critic diverges, reusing a million transitions is worthless.
\end{keyidea}

This is why PPO dominates in RLHF, large-scale simulation, and robotics despite its lower theoretical sample efficiency: stability and parallelizability often matter more than per-sample reuse.

%% ----------------------------------------------------------------
\section{The Compressed Decision Rule}
\label{sec:compressed-rule}

The entire on/off-policy distinction compresses to:

\begin{quote}
If the update formula requires the sample actions to be drawn from the current policy, the algorithm is \textbf{on-policy}. If it can use any historical actions as long as $(s, a, r, s')$ is available, it is \textbf{off-policy}. If it can use actions from a recent old policy but must constrain the divergence between old and new, it is \textbf{near on-policy}.
\end{quote}

The underlying question is always the same: \emph{is the action distribution in the historical data allowed to differ from the distribution of the policy currently being optimized?}

\chapter{Exploration, Exploitation, and Entropy}
\label{ch:exploration}

Every RL agent faces a dilemma: exploit what it currently believes is best, or explore to improve its beliefs. This chapter examines how the dilemma arises from the interaction between noisy value estimates and greedy action selection, and traces the progression from crude fixes ($\varepsilon$-greedy) to principled solutions (maximum entropy RL).

%% ----------------------------------------------------------------
\section{Greedy Converts Noise into Bias}
\label{sec:greedy-problem}

Consider a naive agent that always selects $a = \arg\max_a Q(s,a)$. In early training, $Q$ is poorly estimated---essentially noisy. But $\arg\max$ does not know that. It selects whichever action's estimate happens to be highest, regardless of whether that estimate is accurate.

The resulting feedback loop is destructive:

\begin{equation*}
  \underset{\text{1. early $Q$ noisy}}{\boxed{\phantom{x}}}
  \;\to\;
  \underset{\text{2. argmax picks overestimated action}}{\boxed{\phantom{x}}}
  \;\to\;
  \underset{\text{3. policy biases toward it}}{\boxed{\phantom{x}}}
  \;\to\;
  \underset{\text{4. other actions unsampled}}{\boxed{\phantom{x}}}
  \;\to\;
  \underset{\text{5. data narrows}}{\boxed{\phantom{x}}}
  \;\to\;
  \underset{\text{6. $Q$ harder to correct}}{\boxed{\phantom{x}}}
\end{equation*}

\begin{keyidea}[The Core Problem with Greedy]
The problem with greedy action selection is not simply ``too greedy.'' It is that greedy selection converts estimation noise directly into sampling bias. Once the data distribution narrows, the noisy estimates that caused the problem become self-reinforcing.
\end{keyidea}

From a statistical perspective: without sampling, there is no estimation; without coverage, there is no generalization guarantee.

%% ----------------------------------------------------------------
\section{$\varepsilon$-Greedy: A Hard Rule}
\label{sec:epsilon-greedy}

DQN's standard exploration mechanism is $\varepsilon$-greedy:
\begin{equation}
  a =
  \begin{cases}
    \arg\max_a Q(s,a) & \text{with probability } 1 - \varepsilon, \\
    \text{Uniform}(\mathcal{A}) & \text{with probability } \varepsilon.
  \end{cases}
  \label{eq:eps-greedy}
\end{equation}

This is a binary switch: either full exploitation or fully random exploration, with no gradation. It is crude but effective at preventing the feedback loop of Section~\ref{sec:greedy-problem}. The random actions inject data diversity that keeps Q-estimates from becoming self-fulfilling prophecies.

Typical practice decays $\varepsilon$ over training---exploring widely at the start, then narrowing as Q-estimates improve. The schedule is a hyperparameter, and getting it wrong (decaying too fast or too slowly) directly impacts performance.

\subsection{Limitations}

\begin{itemize}
  \item The random actions are \emph{uniformly} random, with no preference for informative actions.
  \item The binary switch creates a sharp boundary between exploitation and exploration modes.
  \item In high-dimensional action spaces, uniform random actions are almost certainly useless.
\end{itemize}

%% ----------------------------------------------------------------
\section{Entropy: A Soft Constraint on Policy Sharpness}
\label{sec:entropy-exploration}

Entropy provides a continuous, differentiable alternative to the hard switch of $\varepsilon$-greedy. For a stochastic policy $\pi_\theta(a \mid s)$, the entropy is:
\begin{equation}
  \mathcal{H}\bigl[\pi_\theta(\cdot \mid s)\bigr]
  = -\sum_a \pi_\theta(a \mid s)\,\log\pi_\theta(a \mid s)
  \qquad\text{(discrete)},
\end{equation}
or the differential entropy analogue for continuous actions.

When entropy is high, the policy spreads probability across many actions. When entropy is low, the policy is sharp and nearly deterministic. Adding an entropy bonus to the loss encourages the policy to maintain breadth:
\begin{equation}
  \mathcal{L} = -\hat{A}(s,a)\,\log\pi_\theta(a \mid s) - \beta\,\mathcal{H}\bigl[\pi_\theta(\cdot \mid s)\bigr],
\end{equation}
where $\beta > 0$ controls the strength of the regularization.

\begin{keyidea}[Entropy vs.\ $\varepsilon$-Greedy]
$\varepsilon$-greedy is a hard, external rule: either greedy or random. Entropy is a soft, internal constraint: the policy itself learns how much randomness to maintain. Entropy smoothly modulates the exploration--exploitation balance rather than switching between two extremes.
\end{keyidea}

\subsection{Why Not Just ``Encourage Randomness''?}

The entropy bonus is sometimes described as ``encouraging exploration,'' which is true but imprecise. Its deeper role is \emph{distributional regularization}: preventing the policy from collapsing to a near-deterministic distribution before the value estimates have converged. Without it, the advantage-weighted gradient pushes probability mass toward a few actions, narrowing the sampling distribution and degrading future estimates.

%% ----------------------------------------------------------------
\section{SAC: Exploration as Part of the Objective}
\label{sec:sac-exploration}

Soft Actor-Critic (SAC) makes entropy a first-class citizen of the optimization objective rather than an auxiliary regularizer:
\begin{equation}
  J_{\text{SAC}}(\theta) = \sum_t \mathbb{E}\!\bigl[
    r(s_t, a_t) + \alpha\,\mathcal{H}\bigl[\pi_\theta(\cdot \mid s_t)\bigr]
  \bigr],
  \label{eq:sac-objective}
\end{equation}
where $\alpha$ is a temperature parameter that balances reward against entropy.

This reframes the problem: the agent is not maximizing reward and then adding exploration as an afterthought. It is maximizing \emph{entropy-augmented reward}. Sampling diversity is something the agent is explicitly rewarded for.

\subsection{Automatic Temperature Tuning}

SAC can automatically adjust $\alpha$ by treating it as a Lagrange multiplier for a minimum-entropy constraint:
\begin{equation}
  \alpha^* = \arg\min_\alpha\; \mathbb{E}\!\bigl[
    -\alpha\,\log\pi_\theta(a \mid s) - \alpha\,\bar{\mathcal{H}}
  \bigr],
\end{equation}
where $\bar{\mathcal{H}}$ is a target entropy (typically set to $-\dim(\mathcal{A})$ for continuous actions). When the policy's entropy drops below the target, $\alpha$ increases to push it back up. When entropy is already high, $\alpha$ decreases to allow more exploitation.

\subsection{Consequences for the Q-Function}

SAC's critic learns a \emph{soft} Q-function that accounts for future entropy:
\begin{equation}
  Q^{\text{soft}}(s,a) = r + \gamma\,\mathbb{E}_{s'}\!\bigl[
    V^{\text{soft}}(s')
  \bigr],
  \qquad
  V^{\text{soft}}(s) = \mathbb{E}_{a \sim \pi}\!\bigl[
    Q^{\text{soft}}(s,a) - \alpha\log\pi(a \mid s)
  \bigr].
\end{equation}
The soft value includes the entropy bonus, so the critic naturally accounts for the exploration benefit of stochastic actions. This is more principled than bolting entropy onto a standard Q-learner.

%% ----------------------------------------------------------------
\section{The Balance: Preference Shaping vs.\ Uniform Pull}
\label{sec:balance}

Every stochastic policy in RL lives under the tension of two forces:

\begin{center}
\begin{tabular}{@{} l p{8cm} @{}}
\toprule
\textbf{Force} & \textbf{Effect on the policy distribution} \\
\midrule
Preference shaping & Push probability toward actions with high estimated advantage. This is the reward-driven signal: actions that worked well should become more likely. \\
Entropy / uniform pull & Pull the distribution back toward uniformity. This prevents premature commitment and ensures continued data coverage. \\
\bottomrule
\end{tabular}
\end{center}

The final learned policy is the equilibrium between these two forces:
\begin{equation}
  \pi^*(a \mid s) \;\propto\; \exp\!\Bigl(\frac{1}{\alpha}\,Q^{\text{soft}}(s, a)\Bigr).
  \label{eq:boltzmann}
\end{equation}

This is a Boltzmann (softmax) distribution over actions, parameterized by the soft Q-values. At high temperature ($\alpha$ large), the distribution is nearly uniform---maximum exploration. At low temperature ($\alpha$ small), the distribution concentrates on the highest-Q action---maximum exploitation. The optimal $\alpha$ balances the two.

\begin{keyidea}[The Exploration--Exploitation Equilibrium]
The final policy is neither pure exploitation (a delta function on the best action) nor pure exploration (a uniform distribution). It is a temperature-controlled Boltzmann distribution where the temperature $\alpha$ mediates between preference shaping and distributional breadth. SAC automates the choice of $\alpha$; simpler methods like $\varepsilon$-greedy approximate it with a manual schedule.
\end{keyidea}

%% ----------------------------------------------------------------
\section{Summary: The Exploration Progression}
\label{sec:exploration-summary}

\begin{center}
\begin{tabular}{@{} l l l @{}}
\toprule
\textbf{Method} & \textbf{Mechanism} & \textbf{Character} \\
\midrule
Greedy & $\arg\max Q$ & No exploration; noise $\to$ bias \\
$\varepsilon$-greedy & Random with probability $\varepsilon$ & Hard switch; crude but functional \\
Entropy bonus & $+\beta\,\mathcal{H}[\pi]$ in loss & Soft regularizer on policy sharpness \\
SAC (max-entropy) & $\max\; r + \alpha\,\mathcal{H}[\pi]$ & Exploration as part of objective \\
\bottomrule
\end{tabular}
\end{center}

The progression is from external, ad-hoc fixes to internal, principled objectives. At each step, the agent gains a more nuanced mechanism for balancing what it knows (exploit) against what it does not yet know (explore). The endpoint---maximum entropy RL---treats exploration not as a cost to be minimized but as a property to be optimized alongside reward.

\chapter{Bias Management --- RL's Real Game}
\label{ch:bias-management}

Newcomers to RL often assume the goal is to find an unbiased estimator of the
optimal policy.  This is a misunderstanding.  Every component of the RL
pipeline---sampling, estimation, target construction---introduces bias.  The
real engineering question is not ``how do I eliminate bias?'' but rather
\emph{which biases am I willing to accept, and how do I keep them under
control?}

\begin{keyidea}[The Bias Management Principle]
RL does not pursue unbiased optimization.  It selects, among a landscape of
unavoidable biases, the bias structure that is most stable and controllable in
practice.
\end{keyidea}


%% ================================================================
\section{Seven Sources of Bias}
\label{sec:seven-biases}

Almost every design choice in RL introduces a distinct form of bias.
Table~\ref{tab:bias-sources} catalogs the seven most important ones.

\begin{table}[ht]
\centering
\caption{Seven sources of bias in reinforcement learning.}
\label{tab:bias-sources}
\begin{tabular}{@{}lp{8.5cm}@{}}
\toprule
\textbf{Bias source} & \textbf{Mechanism} \\
\midrule
Sampling bias &
  The agent only visits states and actions that its current (or past) policy
  selects.  Regions of the state--action space that are never sampled can never
  be evaluated. \\
Bootstrapping bias &
  TD targets use the model's own estimates as labels:
  $y = r + \gamma \hat{Q}(s',a')$.  Errors in $\hat{Q}$ feed back into the
  next round of fitting. \\
Max / argmax bias &
  $\mathbb{E}[\max_a \hat{Q}(s,a)] \geq \max_a \mathbb{E}[\hat{Q}(s,a)]$.
  Taking the max over noisy estimates systematically overestimates value. \\
Function approximation bias &
  Neural networks generalize---and therefore \emph{mis}-generalize.  A network
  that fits well in visited regions may produce arbitrary values elsewhere. \\
Replay distribution bias &
  Off-policy methods reuse transitions sampled by old policies.  The
  distribution in the buffer drifts away from the current policy's visitation
  distribution. \\
Reward design bias &
  Shaped rewards, sparse rewards, and proxy rewards all inject the designer's
  assumptions.  The agent optimizes the reward it is given, not the reward the
  designer intended. \\
Policy-induced data bias &
  The policy generates its own training data.  A prematurely peaked policy
  narrows the data distribution, making it harder for the estimator to
  self-correct. \\
\bottomrule
\end{tabular}
\end{table}

These biases are not independent---they interact.  Bootstrapping amplifies
function-approximation error; max bias turns that amplified error into
overestimation; overestimation warps the policy, which in turn narrows sampling.
Understanding these feedback loops is the key to understanding why RL algorithms
look the way they do.


%% ================================================================
\section{Q-Learning's Bias: Max Overestimation}
\label{sec:q-bias}

Q-learning constructs its target as
\begin{equation}
  y = r + \gamma \max_{a'} Q_{\text{target}}(s', a').
\end{equation}
If each $\hat{Q}(s',a')$ has zero-mean noise $\epsilon_{a'}$, then
\begin{equation}
  \mathbb{E}\!\bigl[\max_{a'} \hat{Q}(s',a')\bigr]
  \;>\; \max_{a'}\mathbb{E}\!\bigl[\hat{Q}(s',a')\bigr],
\end{equation}
because the max operator preferentially selects actions whose noise happens to
be positive.  This is Jensen's inequality applied to the convex $\max$
function.

Over many bootstrapping steps, the overestimation compounds.  The agent begins
to ``believe'' that certain state--action pairs are far more valuable than they
actually are, and the policy chases these phantoms.

\subsection{Correction 1: Double Q-Learning}

Double DQN decouples action selection from value evaluation:
\begin{equation}
  a^{\ast} = \operatorname{argmax}_{a'} Q_{\text{online}}(s', a'),
  \qquad
  y = r + \gamma\, Q_{\text{target}}(s', a^{\ast}).
\end{equation}
The online network picks the action; the target network evaluates it.  Because
two different networks are unlikely to share the same noise pattern,
overestimation is reduced (though not eliminated).

\subsection{Correction 2: Target Network}

A target network $Q_{\text{target}}$ is a slow-moving copy of $Q_{\text{online}}$.
By not chasing every gradient step, it stabilizes the bootstrapping target and
dampens the positive-feedback loop between overestimation and policy shift.

\subsection{Correction 3: Clipped Double Q / Conservative Q}

TD3 maintains two critics and takes the minimum:
\begin{equation}
  y = r + \gamma\, \min\!\bigl(Q_1(s', \mu_{\text{target}}(s')),\;
  Q_2(s', \mu_{\text{target}}(s'))\bigr).
\end{equation}
This deliberately introduces a \emph{pessimistic} bias to counteract the
optimistic bias of $\max$.  The trade-off is slower value propagation, but
the critic surface the actor exploits is far more trustworthy.

Conservative Q-Learning (CQL) pushes the idea further in the offline setting by
adding a penalty that explicitly lowers $Q$ values for out-of-distribution
actions, ensuring the policy does not exploit regions where the critic has no
data support.

\begin{keyidea}[The Q-Learning Bias Correction Arc]
$\max$ overestimation $\;\longrightarrow\;$
Double Q (decouple selection/evaluation) $\;\longrightarrow\;$
Target net (slow the bootstrap) $\;\longrightarrow\;$
Clipped double Q (pessimistic bias) $\;\longrightarrow\;$
Conservative Q (penalize unsupported actions).
Each step trades a different form of controlled bias for a reduction in
uncontrolled overestimation.
\end{keyidea}


%% ================================================================
\section{Policy Gradient's Bias: Variance That Masquerades as Bias}
\label{sec:pg-bias}

Policy gradient methods avoid the $\max$ operator entirely.  The gradient
\begin{equation}
  \nabla J(\theta) = \mathbb{E}_{s,a \sim \pi_\theta}
  \bigl[\nabla \log \pi_\theta(a|s)\; A^{\pi}(s,a)\bigr]
\end{equation}
is, in principle, \emph{unbiased}---provided the advantage $A^{\pi}$ is computed
exactly and the actions are sampled from the current policy.

In practice, neither condition holds.  The advantage is estimated from a finite
number of noisy rollouts, and Monte Carlo returns have high variance.  High
variance means that the gradient signal is dominated by noise, and training
becomes unstable or painfully slow.

\subsection{Baseline: Centering the Advantage}

Subtracting a baseline $b(s)$ from the return does not change the expected
gradient (the $\nabla\log\pi$ term integrates to zero under the policy
distribution), but it dramatically reduces variance:
\begin{equation}
  \nabla J(\theta) = \mathbb{E}\bigl[\nabla \log \pi_\theta(a|s)\;
  \bigl(A(s,a) - b(s)\bigr)\bigr].
\end{equation}
The value function $V(s)$ is the most common baseline, yielding the advantage
$A(s,a) = Q(s,a) - V(s)$.

\subsection{GAE: Bias--Variance Trade-off in Advantage Estimation}

Generalized Advantage Estimation (GAE) interpolates between low-bias,
high-variance Monte Carlo returns and high-bias, low-variance one-step TD
errors via the parameter $\lambda$:
\begin{equation}
  \hat{A}_t^{\text{GAE}(\gamma,\lambda)}
  = \sum_{l=0}^{\infty} (\gamma\lambda)^l\, \delta_{t+l},
  \qquad
  \delta_t = r_t + \gamma V(s_{t+1}) - V(s_t).
\end{equation}
When $\lambda = 1$, GAE reduces to the full Monte Carlo advantage (unbiased,
high variance).  When $\lambda = 0$, it reduces to the one-step TD advantage
(biased, low variance).  Practitioners typically choose $\lambda \in [0.9, 0.99]$,
accepting a small amount of bias in exchange for much lower variance.

\subsection{PPO Clipping: Trading Bias for Stability}

PPO uses the importance-sampled objective
\begin{equation}
  L^{\text{CLIP}}(\theta) = \hat{\mathbb{E}}\!\left[\min\!\left(
    r_t(\theta)\,\hat{A}_t,\;
    \operatorname{clip}\!\bigl(r_t(\theta),\, 1{-}\epsilon,\, 1{+}\epsilon\bigr)\,
    \hat{A}_t
  \right)\right],
\end{equation}
where $r_t(\theta) = \pi_\theta(a_t|s_t)\,/\,\pi_{\text{old}}(a_t|s_t)$.

The clipping deliberately prevents the ratio from moving too far from 1.  This
means the gradient is no longer the true policy gradient---it is biased.  But
the bias is bounded and predictable, whereas the unclipped gradient can produce
catastrophic policy updates when the importance ratio explodes.

\begin{keyidea}[PPO's Bargain]
PPO intentionally introduces gradient bias through clipping.  The payoff is
engineering stability: updates that are ``close enough'' to the true gradient but
never catastrophically wrong.  This is perhaps the clearest example in RL of
\emph{choosing a more controllable bias structure}.
\end{keyidea}


%% ================================================================
\section{SAC's Bias: Entropy as Intentional Objective Rewriting}
\label{sec:sac-bias}

Soft Actor-Critic optimizes a modified objective:
\begin{equation}
  J_{\text{SAC}}(\pi) = \sum_t \mathbb{E}\!\bigl[
    r(s_t, a_t) + \alpha\,\mathcal{H}\!\bigl(\pi(\cdot|s_t)\bigr)
  \bigr],
\end{equation}
where $\alpha$ is a temperature parameter and $\mathcal{H}$ is the entropy of
the policy.

This is \emph{not} the original reward-maximization objective.  SAC admits an
intentional bias: the optimal SAC policy is not the same as the optimal
reward-maximizing policy.  It is the optimal \emph{entropy-regularized} policy.

Why accept this bias?  Three reasons:
\begin{enumerate}[nosep]
  \item \textbf{Stable exploration.}  The entropy term prevents the policy from
    collapsing onto a single action.  This keeps the sampling distribution broad
    and the critic's training data diverse.
  \item \textbf{Smoother Q landscape.}  When the policy stays stochastic, the
    Bellman target integrates over actions rather than hinging on a single
    $\max$.  This suppresses the spike-chasing behavior that plagues
    deterministic actor-critics.
  \item \textbf{Robustness.}  A stochastic policy is inherently more robust to
    perturbations---it hedges rather than committing to a single action.
\end{enumerate}

SAC's bias structure can be summarized as:

\begin{keyidea}[SAC's Entropy Bias]
SAC trades a soft objective (entropy-regularized reward) for stable exploration
and better data coverage.  The bias is explicit, tunable via $\alpha$, and
produces policies that are both effective and robust.
\end{keyidea}


%% ================================================================
\section{Choosing a Controllable Bias Structure}
\label{sec:choosing-bias}

Table~\ref{tab:bias-strategies} summarizes how major algorithm families handle
the bias problem.

\begin{table}[ht]
\centering
\caption{Bias strategies across algorithm families.}
\label{tab:bias-strategies}
\begin{tabular}{@{}lp{4cm}p{4.5cm}@{}}
\toprule
\textbf{Algorithm family} & \textbf{Dominant bias} & \textbf{Mitigation strategy} \\
\midrule
Q-learning / DQN &
  Max overestimation, bootstrap &
  Double Q, target net, conservative Q \\
DDPG / TD3 &
  Actor exploits critic errors &
  Clipped double critic, delayed update, target smoothing \\
REINFORCE / A2C &
  High variance (effective bias in finite samples) &
  Baseline, GAE, advantage normalization \\
PPO &
  Clipping introduces gradient bias &
  Bounded, predictable bias traded for stability \\
SAC &
  Entropy rewrites the objective &
  Intentional bias for exploration; tunable via $\alpha$ \\
\bottomrule
\end{tabular}
\end{table}

The pattern is clear.  No algorithm eliminates bias.  Each algorithm
\emph{chooses} which biases to accept based on what is controllable in its
particular setting:

\begin{itemize}[nosep]
  \item Value-based methods accept bootstrapping and approximation bias, then
    spend engineering effort suppressing max overestimation.
  \item Policy gradient methods are unbiased in theory but high-variance in
    practice; they introduce controlled bias (baselines, clipping, GAE's
    $\lambda$) to tame variance.
  \item SAC openly rewrites the objective, accepting an entropy bias that is
    explicit and tunable.
\end{itemize}

\begin{keyidea}[The Meta-Lesson]
The history of RL algorithm design is not a march toward unbiased estimation.
It is a search for \textbf{better bias structures}---ones whose distortions are
predictable, bounded, and tunable, even if the resulting objective is no longer
the ``pure'' reward-maximization problem.
\end{keyidea}


%% ================================================================
\section{Why This Matters in Practice}
\label{sec:bias-practice}

When debugging an RL system, the first question should not be ``is my
implementation correct?'' but rather ``which bias is dominating, and is it the
one I designed for?''

A few practical heuristics:

\begin{enumerate}[nosep]
  \item \textbf{If Q values diverge upward}: max overestimation is likely
    compounding.  Switch to Double Q or clipped double critics.
  \item \textbf{If the policy collapses to a single action}: sampling bias has
    taken over.  Increase entropy regularization or exploration noise.
  \item \textbf{If training is stable but performance plateaus}: the bias
    structure may be too conservative.  Try reducing PPO's clip range, lowering
    SAC's $\alpha$, or using less aggressive target-network lag.
  \item \textbf{If off-policy training is unstable}: replay distribution bias
    may be severe.  Consider prioritized replay, shorter buffer horizons, or
    switching to a near-on-policy method.
\end{enumerate}

Understanding bias is not an academic exercise---it is the most practical skill
in RL engineering.  Every hyperparameter in a modern RL algorithm is, at its
core, a knob that adjusts the bias--variance--stability trade-off.

\chapter{Q as Value Compression, Policy as Preference Compression}
\label{ch:compression}

Most treatments of RL describe the Q function and the policy as two different
``things to learn.''  This chapter offers a more illuminating framing: both are
\textbf{compressors}, but they compress different objects.  Understanding what
each one compresses---and how the compressed representations connect---explains
why algorithms are structured the way they are and why the field moved from
value-first methods to policy-first methods.


%% ================================================================
\section{Q: Compressing Sparse Rewards into Dense Value}
\label{sec:q-compression}

Consider a typical RL environment.  The agent receives a reward of $+1$ only
when it reaches a distant goal; every other timestep yields $r=0$.  Raw rewards
are \emph{sparse}---they carry almost no per-step information.

The Q function compresses the entire future into a single scalar:
\begin{equation}
  Q(s,a) = \mathbb{E}\!\left[\sum_{k=0}^{\infty}\gamma^k\, r_{t+k}
  \;\Big|\; s_t = s,\; a_t = a\right].
\end{equation}

This is a lossy compression.  The full distribution of possible futures---all
the branching trajectories, stochastic transitions, and delayed rewards---is
collapsed into one number.  But it is an extraordinarily useful compression:
it tells you, at every state--action pair, how much long-term value to expect.

\begin{keyidea}[Q as Value Compressor]
The Q function compresses sparse rewards and complex future rollouts into a
dense, per-step scalar signal.  It transforms a credit-assignment problem
(``which past action caused this delayed reward?'') into a local lookup
(``what is this action worth right here?'').
\end{keyidea}


%% ================================================================
\section{Policy: Compressing Behavioral Preferences}
\label{sec:policy-compression}

The policy is also a compressor, but its target is different.  Rather than
compressing \emph{value}, it compresses \emph{preference}---the agent's
behavioral tendency in each state.

A stochastic policy $\pi(a|s)$ represents preferences as a probability
distribution.  A deterministic policy $\mu(s)$ represents preferences as a
single point.  Both are compressed summaries of ``how I prefer to act here,''
distilled from experience.

\begin{keyidea}[Policy as Preference Compressor]
The policy compresses the agent's accumulated behavioral preferences into a
compact representation---a distribution $\pi(a|s)$ or a point $\mu(s)$---that
can be evaluated in a single forward pass.
\end{keyidea}


%% ================================================================
\section{Two Compression Strategies}
\label{sec:two-strategies}

The distinction between value compression and preference compression maps
directly onto the two major families of RL algorithms.

\subsection{Value-First: Compress Value, Then Extract Preference}

Q-learning and DQN follow a two-stage strategy:
\begin{enumerate}[nosep]
  \item \textbf{Compress value}: learn $Q(s,a)$ for all actions.
  \item \textbf{Extract preference}: apply $\operatorname{argmax}$ to convert
    the value landscape into a behavioral choice.
\end{enumerate}

\begin{equation}
  \underbrace{\text{Sparse rewards} \;\longrightarrow\; Q(s,a)}_{\text{value compression}}
  \;\longrightarrow\;
  \underbrace{a^{\ast} = \operatorname{argmax}_a Q(s,a)}_{\text{preference extraction}}
\end{equation}

The advantage is that $Q$ can be learned from off-policy data with high sample
efficiency.  The disadvantage is that the second stage---argmax---is a fragile
bridge between value and preference (Section~\ref{sec:fragile-argmax}).

\subsection{Policy-First: Compress Preference Directly}

Policy gradient methods skip the value-compression stage (or use it only as an
auxiliary signal) and directly optimize the preference representation:
\begin{equation}
  \underbrace{\text{Sampled returns / advantages}
  \;\longrightarrow\; \pi_\theta(a|s)}_{\text{preference compression}}
\end{equation}

The advantage is robustness: there is no brittle argmax step.  The disadvantage
is data hunger---on-policy methods must discard data whenever the policy
changes, and the gradient signal from sampled returns is noisy.

\subsection{Actor-Critic: Value Compressor Guides Preference Compressor}

Actor-Critic architectures maintain \emph{both} compressors simultaneously:
\begin{itemize}[nosep]
  \item The \textbf{critic} compresses value (or advantage).
  \item The \textbf{actor} compresses preference.
  \item The critic's compressed value signal \emph{guides} the actor's
    preference updates.
\end{itemize}

\begin{equation}
  \underbrace{Q_\phi(s,a) \text{ or } A_\phi(s,a)}_{\text{value compressor (critic)}}
  \;\xrightarrow{\text{guides}}\;
  \underbrace{\pi_\theta(a|s)}_{\text{preference compressor (actor)}}
\end{equation}

This is not simply ``one network scores, another network acts.''  It is a
cooperative compression scheme in which the value compressor provides a
low-variance training signal that accelerates the preference compressor's
learning.

The risk is mutual contamination: if the critic's value compression is poor
(overestimation, dead regions), it corrupts the actor's preference compression.
This is why Actor-Critic methods need careful stabilization---target networks,
double critics, entropy bonuses---to keep both compressors healthy.


%% ================================================================
\section{Why Hard Argmax Is a Fragile Bridge}
\label{sec:fragile-argmax}

The argmax operation is the standard way to convert a Q-value landscape into a
decision:
\begin{equation}
  a^{\ast} = \operatorname{argmax}_a Q(s,a).
\end{equation}

In theory, if $Q = Q^{\ast}$, this recovers the optimal policy.  In practice,
$Q$ is an approximation with noise, and argmax has two pathological properties:

\begin{enumerate}
  \item \textbf{Discontinuity.}  Argmax is a discontinuous function of $Q$.  A
    tiny perturbation in Q-values can flip the selected action entirely.  This
    means that small estimation errors produce large behavioral changes.
  \item \textbf{Noise amplification.}  Among $k$ noisy estimates, argmax
    preferentially selects the one with the largest positive noise.  This is the
    overestimation bias discussed in Chapter~\ref{ch:bias-management}, but seen
    here as a \emph{compression failure}: the value compressor's noise is
    amplified at the value-to-preference interface.
\end{enumerate}

A soft policy $\pi(a|s) \propto \exp(Q(s,a)/\alpha)$ replaces the hard argmax
with a smooth, temperature-controlled mapping.  Small errors in $Q$ produce
small changes in the distribution, not discontinuous action switches.  This is
one of the deepest reasons the field moved toward stochastic, soft policies.

\begin{keyidea}[Argmax as Compression Bottleneck]
Hard argmax is a fragile bridge between value compression and preference
compression.  It amplifies noise and is discontinuous.  Soft policies replace
this brittle interface with a smooth, differentiable mapping---making the
overall compression pipeline more robust.
\end{keyidea}


%% ================================================================
\section{Algorithm Comparison Through the Compression Lens}
\label{sec:compression-table}

Table~\ref{tab:compression} summarizes major algorithms in terms of what they
compress, how they make decisions, and what can go wrong.

\begin{table}[ht]
\centering
\caption{Algorithms viewed as compression systems.}
\label{tab:compression}
\begin{tabular}{@{}llll@{}}
\toprule
\textbf{Algorithm} & \textbf{Compression object} &
\textbf{Decision method} & \textbf{Core risk} \\
\midrule
Q-learning / DQN &
  Long-term value $Q$ &
  Hard argmax &
  Q noise amplified by argmax \\
Double DQN / TD3 &
  Value (stabilized) &
  Conservative argmax &
  Underestimation; slower learning \\
Policy Gradient &
  Preference $\pi$ directly &
  Sample from $\pi$ &
  High variance; data hungry \\
PPO &
  Preference $\pi$ (stabilized) &
  Clipped update of $\pi$ &
  May be overly conservative \\
SAC &
  Soft preference &
  Sample from entropy-regularized $\pi$ &
  Objective rewritten by entropy \\
Actor-Critic &
  $Q/A$ (value) + $\pi$ (preference) &
  Critic guides actor &
  Compressors can contaminate each other \\
\bottomrule
\end{tabular}
\end{table}


%% ================================================================
\section{The Compression Perspective in One Sentence}
\label{sec:compression-summary}

\begin{keyidea}[The Compression View of RL]
$Q$ is the compression of sparse rewards and long-term value into a dense
scalar.  $\pi$ is the compression of behavioral preferences into a
distribution.  The history of RL algorithm design is the story of learning how
to build, connect, and stabilize these two compressors---and of discovering
that directly compressing preferences is often more robust than compressing
value first and then extracting preferences through the fragile bottleneck of
argmax.
\end{keyidea}

\chapter{Action Modeling}
\label{ch:action-modeling}

Everything discussed so far---value compression, preference compression, bias
management---ultimately feeds into one concrete question: \emph{how does the
agent choose an action?}  The answer depends on whether the action space is
discrete or continuous, and whether the policy is deterministic or stochastic.
These choices are not cosmetic; they fundamentally shape the algorithm's
architecture, exploration behavior, and failure modes.


%% ================================================================
\section{Discrete Actions: Enumerate, Then Argmax}
\label{sec:discrete}

When the action space is a finite set $\mathcal{A} = \{a_1, a_2, \ldots, a_k\}$,
the Q-learning approach is natural.  A neural network takes the state as input
and outputs one Q-value per action:
\begin{equation}
  s \;\longrightarrow\; \text{network} \;\longrightarrow\;
  \bigl[Q(s,a_1),\; Q(s,a_2),\; \ldots,\; Q(s,a_k)\bigr].
\end{equation}
Decision-making is trivial: scan the $k$ outputs and pick the largest.
\begin{equation}
  a^{\ast} = \operatorname{argmax}_{a \in \mathcal{A}} Q(s,a).
\end{equation}
DQN is built on exactly this structure.  The convenience of discrete actions is
that argmax is a free operation---$O(k)$ comparisons with no optimization
required.  The limitation is equally clear: the approach does not scale when
$k$ is very large or when actions are continuous.


%% ================================================================
\section{Continuous Actions: The Need for an Actor}
\label{sec:continuous}

In continuous control, actions are real-valued vectors $a \in \mathbb{R}^n$.
Argmax becomes a continuous optimization problem:
\begin{equation}
  a^{\ast} = \operatorname{argmax}_{a \in \mathbb{R}^n} Q(s, a),
\end{equation}
which must be solved at every decision step---generally infeasible in real time.

The solution is to introduce an \textbf{actor network} $\mu_\theta(s)$ that
directly outputs an action, effectively serving as a learned, differentiable
approximation to argmax:
\begin{equation}
  a = \mu_\theta(s), \qquad
  \theta^{\ast} = \operatorname{argmax}_\theta\; Q(s,\, \mu_\theta(s)).
\end{equation}
The actor is trained by backpropagating through the critic's gradient
$\nabla_a Q(s,a)\big|_{a=\mu_\theta(s)}$, pushing its output toward regions
where $Q$ is high.

\begin{keyidea}[The Actor as Differentiable Argmax]
In continuous action spaces, the actor network replaces the intractable
$\operatorname{argmax}_{a \in \mathbb{R}^n} Q(s,a)$ with a learned function
$\mu_\theta(s)$ trained via gradient ascent on $Q$.  DDPG is the canonical
realization of this idea.
\end{keyidea}


%% ================================================================
\section{Policy Representations}
\label{sec:policy-representations}

When the algorithm optimizes a policy directly (as in policy gradient, PPO, or
SAC), the policy must be represented as a parameterized distribution from which
actions can be sampled and whose log-probability can be differentiated.

\subsection{Discrete Policy: Categorical Distribution}

For discrete actions, the standard representation is a categorical distribution
produced by a softmax over network logits:
\begin{equation}
  s \;\longrightarrow\; \text{network} \;\longrightarrow\;
  \ell_1, \ldots, \ell_k
  \;\longrightarrow\;
  \pi(a_i|s) = \frac{\exp(\ell_i)}{\sum_j \exp(\ell_j)}.
\end{equation}
Sampling draws an action from this categorical distribution.  The log-probability
$\log \pi(a|s) = \ell_a - \log\sum_j \exp(\ell_j)$ is smooth and differentiable
with respect to the network parameters.

\subsection{Continuous Policy: Gaussian Distribution}

For continuous actions, the most common choice is a diagonal Gaussian:
\begin{equation}
  \pi(a|s) = \mathcal{N}\!\bigl(a;\; \mu_\theta(s),\; \sigma_\theta(s)^2 I\bigr),
\end{equation}
where the network outputs both the mean $\mu_\theta(s)$ and the standard
deviation $\sigma_\theta(s)$ (or log-std for numerical stability):
\begin{equation}
  s \;\longrightarrow\; \text{network} \;\longrightarrow\;
  (\mu,\; \log\sigma)
  \;\longrightarrow\;
  a = \mu + \sigma \odot \epsilon,
  \quad \epsilon \sim \mathcal{N}(0, I).
\end{equation}
The mean encodes the center of the agent's preference; the standard deviation
encodes the exploration range.  This is the policy representation used by PPO
and SAC in continuous-control tasks.


%% ================================================================
\section{Deterministic vs.\ Stochastic Policies}
\label{sec:det-vs-stoch}

\subsection{Deterministic: Hard Preference}

A deterministic policy $a = \mu_\theta(s)$ outputs a single action per state.
It represents a \emph{hard preference}---full commitment to one choice.

\begin{itemize}[nosep]
  \item \textbf{Strengths}: stable execution, simple inference, no sampling
    variance.
  \item \textbf{Weaknesses}: no intrinsic exploration.  The agent must rely on
    external noise (e.g., Ornstein--Uhlenbeck in DDPG, Gaussian noise in TD3)
    to visit new state--action pairs.
  \item \textbf{Representatives}: DDPG, TD3.
\end{itemize}

\subsection{Stochastic: Soft Preference}

A stochastic policy $a \sim \pi_\theta(a|s)$ outputs a distribution.  It
represents a \emph{soft preference}---a graded ranking of actions rather than a
single commitment.

\begin{itemize}[nosep]
  \item \textbf{Strengths}: built-in exploration, compatible with entropy
    regularization, provides a natural log-probability for policy gradient.
  \item \textbf{Weaknesses}: sampling introduces variance into both training
    and execution.
  \item \textbf{Representatives}: PPO, SAC, REINFORCE.
\end{itemize}

Table~\ref{tab:det-stoch} summarizes the comparison.

\begin{table}[ht]
\centering
\caption{Deterministic vs.\ stochastic policies.}
\label{tab:det-stoch}
\begin{tabular}{@{}lll@{}}
\toprule
& \textbf{Deterministic} & \textbf{Stochastic} \\
\midrule
Output       & Single action $\mu(s)$ & Distribution $\pi(a|s)$ \\
Preference   & Hard (point estimate)  & Soft (distribution) \\
Exploration  & External noise required & Intrinsic \\
Gradient signal & $\nabla_a Q$ through actor & $\nabla\log\pi \cdot A$ \\
Typical algorithms & DDPG, TD3 & PPO, SAC, REINFORCE \\
\bottomrule
\end{tabular}
\end{table}


%% ================================================================
\section{Why RL Moved Toward Stochastic Policies}
\label{sec:why-stochastic}

In classical control, deterministic policies are the norm: given a state, output
the optimal action.  Why did RL increasingly favor stochastic policies?

The answer lies in RL's unique structure: the policy serves \textbf{two roles
simultaneously}.

\begin{enumerate}
  \item \textbf{Decision-maker}: choosing what to do right now.
  \item \textbf{Data sampler}: determining what the agent will observe in the
    future.
\end{enumerate}

A deterministic policy is a competent decision-maker but a poor data sampler.
It concentrates all future data around a single trajectory, starving the
learning algorithm of the diversity it needs to improve.

A stochastic policy balances both roles.  It makes good (on average) decisions
\emph{and} maintains the data coverage needed for continued learning.

\begin{keyidea}[The Dual Role of Policy]
In RL, the policy is both the decision rule and the data-generation mechanism.
A stochastic policy is a soft preference representation that naturally balances
exploitation (good decisions) with exploration (diverse data).  This dual role
is the fundamental reason RL moved from deterministic to stochastic policies.
\end{keyidea}

From the compression perspective of Chapter~\ref{ch:compression}, a stochastic
policy is a richer, more complete compression of preferences.  Preferences are
not always a single point; they can---and often should---be a distribution that
reflects the agent's uncertainty about which action is truly best.
\begin{align}
  \text{Deterministic policy} &= \text{hard preference (point compression)}, \\
  \text{Stochastic policy}    &= \text{soft preference (distributional compression)}.
\end{align}
The shift from deterministic to stochastic is, at its core, a shift from lossy
point compression to richer distributional compression of behavioral preferences.


\part{Beyond Classical RL}
\chapter{Optimal Control and RL --- Same Problem, Different Trade-offs}
\label{ch:oc-vs-rl}

Optimal Control (OC) and Reinforcement Learning (RL) are taught in different departments, published in different journals, and spoken about in different jargon.
Yet they solve the same problem: \emph{given the current state, choose an action that maximises long-term cumulative reward}.
The differences are engineering trade-offs along three dimensions---rollout, optimisation, and data---each with its own hidden costs.

%------------------------------------------------------------------
\section{Three Dimensions of Comparison}
\label{sec:oc-rl-3d}

\begin{keyidea}[One Problem --- Three Trade-offs]
Every sequential decision system must (i) \emph{imagine} the future (rollout), (ii) \emph{decide} what to do (optimisation), and (iii) \emph{fuel} the process with information (data).
OC and RL make opposite bets on each axis.
\end{keyidea}

\subsection{Rollout: How to Imagine the Future}

\textbf{Optimal Control} maintains an \emph{explicit} dynamics model
$s_{t+1} = f(s_t, a_t)$ (or a stochastic variant $s_{t+1}\sim P(\cdot|s_t,a_t)$)
and simulates forward online at every control step.
Model Predictive Control (MPC) is the canonical example: at each timestep it rolls out the model over a horizon $T$, solves for the optimal action sequence, executes the first action, and repeats.
The model is typically low-dimensional and hand-calibrated; real-time sensor feedback corrects drift.

\textbf{RL} has no explicit model.\footnote{Model-based RL learns one, blurring the line---see Section~\ref{sec:fusion}.}
Instead, it compresses rollout knowledge \emph{offline} into network parameters through massive environment interaction.
At deployment, a single forward pass through $Q_\theta(s,a)$ or $\pi_\theta(a|s)$ replaces the explicit rollout.
The ``imagination'' is implicit---baked into weights, not equations.

\begin{table}[htbp]
\centering
\caption{Rollout comparison.}
\begin{tabular}{@{}lll@{}}
\toprule
& \textbf{Optimal Control} & \textbf{RL} \\
\midrule
Model & Explicit $f(s,a)$ & Implicit (encoded in network) \\
When & Online, every step & Offline training; online inference \\
Input & Precise low-dim state & Accepts high-dim, fuzzy inputs \\
Correction & Real-time sensor feedback & Generalisation ability \\
\bottomrule
\end{tabular}
\end{table}

\subsection{Optimisation: When Does It Happen?}

OC solves an optimisation problem \emph{online} at every timestep:
\[
  a^\star = \arg\min_{a_{0:T}} \sum_{t=0}^{T} \ell(s_t, a_t)
  \quad\text{subject to } s_{t+1}=f(s_t,a_t).
\]
MPC iterates this at control frequency, making deployment compute \emph{high}.

RL pushes all heavy optimisation into offline training.
Online deployment is a single forward pass---$a = \pi_\theta(s)$---making deployment compute \emph{low}.
The price is that adapting to a new scenario requires retraining or fine-tuning, whereas MPC re-optimises automatically given a new state.

\subsection{Data: Few Parameters vs.\ Massive Interactions}

OC's ``data'' is a small set of calibrated parameters---masses, friction coefficients, gear ratios---typically tens to hundreds of scalars obtained through physical measurement or system identification.
Online data is real-time sensor feedback used for correction, not for training.

RL requires millions to billions of $(s,a,r,s')$ transitions, collected in simulation or the real world.
Training compresses this experience into network parameters.

\begin{keyidea}[RL Model = Compressed Offline Experience]
An RL policy network is a \emph{compressed representation} of offline interaction experience, designed for fast online inference.
At deployment, no re-interaction with the environment is needed---the network has already encoded what to do.
\end{keyidea}

%------------------------------------------------------------------
\section{The Hidden Constraints of Optimal Control}
\label{sec:oc-constraints}

OC looks elegant on paper but works because it makes very strong assumptions.
These assumptions break in complex, high-dimensional, or poorly structured problems.

\subsection{Constraint 1: Explicit Reference and Convex Structure}

MPC and LQR require an explicit reference trajectory and constraints expressible in tractable form (linear or convex):
\[
  J = \sum_{t} (x_t - x_{\text{ref}})^\top Q\,(x_t - x_{\text{ref}}) + u_t^\top R\, u_t,
  \qquad A x \le b.
\]
If constraints are nonlinear or non-convex, or if the reference itself is hard to define, the problem becomes intractable.

Worse, the reference requirement \emph{cascades}:
\[
  \text{Perception} \;\longrightarrow\; \text{Planner} \;\longrightarrow\; \text{Controller}
\]
The controller needs a reference from the planner; the planner needs maps, goals, and semantic information from perception.
Each layer has explicit interface requirements, and errors at any layer propagate and amplify downstream.
System complexity is the \emph{product} of each layer's complexity, not the sum.

RL (especially end-to-end RL) aims to skip this cascade:
raw input $\longrightarrow$ neural network $\longrightarrow$ action.
The cost is more data and training, but it avoids the layered explicit-interface problem.

\subsection{Constraint 2: DOF Explosion}

OC's online solving complexity grows steeply with degrees of freedom.
For state dimension $n$, control dimension $m$, and prediction horizon $T$, MPC's optimisation variables number approximately $(n+m)\times T$.
Standard QP solvers have complexity
\[
  \mathcal{O}\!\bigl(((n+m)T)^3\bigr).
\]
For high-DOF robots, multi-agent scenarios, or pixel-level control, online solving simply cannot keep up.

RL is not bound by this: offline training can take as long as needed, and online inference scales linearly with input dimension (a single matrix multiply per layer).

\subsection{Constraint 3: Poor Parallelisability}

MPC's iterative solvers (QP, ADMM, interior point) are inherently sequential---each iteration depends on the previous one.
They exploit GPUs poorly.

RL training is naturally parallel: vectorised environments roll out simultaneously, samples within a mini-batch are independent (GPU matrix operations), and batch inference handles many inputs at once.

\begin{table}[htbp]
\centering
\caption{Hidden costs summary.}
\begin{tabular}{@{}lcc@{}}
\toprule
\textbf{Constraint} & \textbf{OC} & \textbf{RL} \\
\midrule
Reference / convexity & Required & Not needed \\
High DOF & $\mathcal{O}(((n{+}m)T)^3)$ & Linear inference \\
GPU parallelism & Poor & Natural \\
Massive offline data & Not needed & Required \\
Training stability & N/A & Hard \\
Online adaptation & Strong & Weak \\
\bottomrule
\end{tabular}
\end{table}

%------------------------------------------------------------------
\section{Where They Fuse}
\label{sec:fusion}

In practice, OC and RL are not opponents but partners.

\subsection{Model-Based RL}

Algorithms such as Dreamer and MBPO \emph{learn} a dynamics model from data and use it for rollout inside RL training---combining OC's explicit imagination with RL's data-driven flexibility.

\subsection{Hierarchical Stacks}

A common architecture in autonomous driving and robotics:
\begin{itemize}
  \item \textbf{High level}: RL-trained planner selects waypoints or sub-goals.
  \item \textbf{Low level}: MPC or PID controller tracks the reference with real-time feedback.
\end{itemize}
This lets each layer play to its strengths: RL handles the fuzzy, high-dimensional perception-to-planning mapping; OC handles the precise, low-dimensional tracking.

\subsection{Sim2Real}

Train an RL policy offline in simulation (where interactions are cheap and parallelisable), then deploy it on real hardware with sensor-based correction---blending RL's offline optimisation with OC's online correction philosophy.

%------------------------------------------------------------------
\section{The DP $\to$ RL Relaxation Chain}
\label{sec:dp-relaxation}

Dynamic Programming (DP) is the theoretical ancestor of both OC and RL.
RL can be understood as DP with five successive relaxations of assumptions.

\begin{keyidea}[Five Relaxations from DP to Modern RL]
Each relaxation trades an idealised assumption for a practical mechanism, introducing a new source of bias that subsequent techniques must manage.
\end{keyidea}

\subsection{Relaxation 1: Unknown Model $\to$ Sampled Bellman Backup}

DP assumes known $P(s'|s,a)$ and $R(s,a)$.
In practice these are unknown; we can only observe samples $(s,a,r,s')$.
The exact Bellman expectation becomes a sampled estimate---the origin of \textbf{TD learning}.
\[
  V(s) \leftarrow V(s) + \alpha\bigl[r + \gamma V(s') - V(s)\bigr]
\]

\subsection{Relaxation 2: Huge State Space $\to$ Function Approximation}

DP can sweep every state in a table.
Real state spaces are enormous or continuous.
We replace tables with parameterised functions:
$Q(s,a)\approx Q_\theta(s,a)$, $\pi(a|s)\approx\pi_\theta(a|s)$.
The cost: \emph{function approximation bias}.

\subsection{Relaxation 3: Policy-Dependent Data $\to$ Off-Policy / Replay}

DP operates on the model directly---no ``data distribution'' problem.
In RL, data comes from rollouts of the current or past policy, creating a closed loop where the policy shapes the data that trains it.
Off-policy methods and replay buffers partially decouple this loop.

\subsection{Relaxation 4: Noisy Q $\to$ Stabilised Greedy}

DP's $\arg\max$ over true $Q$ is safe.
RL's $\arg\max$ over a noisy $\hat{Q}$ amplifies estimation errors.
This necessitates $\varepsilon$-greedy, Double Q, target networks, and entropy regularisation---all forms of \emph{stabilised approximate greedy}.

\subsection{Relaxation 5: Hard Greedy $\to$ Soft Preference Optimisation}

Even stabilised greedy remains fragile.
Policy gradient, PPO, and SAC abandon hard $\arg\max$ entirely, instead \emph{directly optimising the policy distribution}.
The shift is from ``solve for the optimal value, then extract the policy'' to ``directly shape the preference distribution through sampling and feedback.''

\begin{table}[htbp]
\centering
\caption{The five relaxations from DP to RL.}
\begin{tabular}{@{}llll@{}}
\toprule
\textbf{Stage} & \textbf{Assumption relaxed} & \textbf{Mechanism introduced} & \textbf{New bias} \\
\midrule
DP $\to$ Tabular RL & Known model & TD / sampled backup & Sampling noise \\
Tabular $\to$ Deep RL & Finite states & Neural function approx. & Approx.\ bias \\
Fixed data $\to$ RL & i.i.d.\ data & Replay / off-policy & Distribution shift \\
True $Q$ $\to$ $\hat{Q}$ & Exact values & Double Q, target net & Overestimation \\
Greedy $\to$ Soft & Hard $\arg\max$ & Entropy, PPO clip & Conservative bias \\
\bottomrule
\end{tabular}
\end{table}

%------------------------------------------------------------------
\section{Summary}

Optimal Control and RL are different engineering realisations of the same sequential decision problem:

\begin{itemize}
  \item \textbf{Know dynamics + low-dim + need real-time adaptation} $\to$ Optimal Control.
  \item \textbf{Unknown dynamics + high-dim + can train offline} $\to$ RL.
  \item \textbf{In practice}, the two fuse: model-based RL, hierarchical stacks, and Sim2Real loops.
\end{itemize}

The DP $\to$ RL relaxation chain shows that every RL technique---TD, replay, target networks, Double Q, entropy, clipping---is a direct response to relaxing one of DP's god's-eye-view assumptions.
Understanding \emph{which} assumption each technique addresses is the key to reading any RL algorithm with clarity.

\chapter{RL Meets the Real World}
\label{ch:rl-real-world}

Reinforcement learning does not exist in a vacuum.
It can be trained from scratch, layered on top of supervised-learning base models, woven into diffusion and flow-matching generators, and shuttled between simulation and reality.
This chapter maps the landscape of how RL connects to other paradigms---and where each connection matters most.

%------------------------------------------------------------------
\section{Pure RL from Scratch}
\label{sec:pure-rl}

The classic setting: no pretrained model, random initialisation, learning purely from environment interaction.

\subsection{Value-Based Methods}

The model learns a scalar estimate of long-term return for every state--action pair:
\[
  Q(s,a) \;\approx\; \mathbb{E}\!\Bigl[\sum_{k=0}^{\infty}\gamma^k r_{t+k} \;\Big|\; s_t=s,\;a_t=a\Bigr].
\]
Decisions follow $a = \arg\max_a Q(s,a)$.
Representatives: Q-learning, DQN, Double DQN, Dueling DQN.

Essence: \emph{learn valuations first, then extract behaviour from valuations}.

\subsection{Policy-Based Methods}

The model directly learns a preference distribution over actions:
\[
  \pi_\theta(a|s) \;\to\; \text{behavioural preference distribution}.
\]
Decisions follow $a \sim \pi_\theta(\cdot|s)$.
Representatives: REINFORCE, PPO, SAC.

Essence: \emph{directly learn the preference distribution}, skipping the intermediate value landscape.

\subsection{Actor-Critic: Combining Both}

Most modern algorithms are Actor-Critic:
\begin{itemize}
  \item \textbf{Critic}: learns value ($V$ or $Q$), provides a preference signal (advantage).
  \item \textbf{Actor}: learns the behaviour distribution, guided by the Critic.
\end{itemize}
Representatives: A2C, DDPG, TD3, SAC, PPO (with a $V$-critic).

%------------------------------------------------------------------
\section{Architecture: Training vs.\ Inference Modules}
\label{sec:architecture}

A recurring pattern across all RL architectures: \emph{training requires more modules than inference}.

\begin{keyidea}[Critic = Scaffolding]
The Critic (or Value Head) exists only during training.
It provides gradient signals---advantage estimates or $Q$-gradients---that guide the Actor's updates.
At deployment, the Critic is discarded; only the Actor/Policy network is shipped to production.
\end{keyidea}

\begin{table}[htbp]
\centering
\caption{What is used at training vs.\ inference.}
\begin{tabular}{@{}llll@{}}
\toprule
\textbf{Architecture} & \textbf{Training} & \textbf{Inference} & \textbf{Discarded} \\
\midrule
DQN & Backbone + Q Head + Target Net & Backbone + Q Head & Target Net \\
PPO (shared) & Backbone + Policy Head + Value Head & Backbone + Policy Head & Value Head \\
SAC / TD3 & Actor Net + Critic Net(s) & Actor Net & Entire Critic \\
\bottomrule
\end{tabular}
\end{table}

%------------------------------------------------------------------
\section{Reward Sources: Prepared Before Training}
\label{sec:reward-sources}

The RL loop requires reward signals, but reward is \emph{not learned during RL training}---it must be prepared beforehand.

\begin{enumerate}
  \item \textbf{Environment-provided reward.}
    Game scores, distance, collision penalties, energy consumption.
    The environment directly returns a scalar; no extra preparation needed.

  \item \textbf{Rule-based reward.}
    Math-answer correctness, code compilation pass rate, format compliance.
    A hand-written function $R(\text{output})\to\mathbb{R}$.
    Common with GRPO.

  \item \textbf{Learned Reward Model (RM).}
    Collect human preference pairs (chosen, rejected), train an SL model
    $\text{RM}(x,y)\to\mathbb{R}$ to approximate human preferences, then freeze it.
    Standard approach for PPO-RLHF (InstructGPT).
\end{enumerate}

\begin{keyidea}[RM Is a Precondition --- Not Part of RL]
The Reward Model is trained via supervised learning on preference data and frozen as a ``scorer'' during RL training.
It is logically \emph{upstream} of RL, not inside it.
\end{keyidea}

%------------------------------------------------------------------
\section{Value and Policy: Coupled in Sampling, Decoupled in Gradients}
\label{sec:coupled-decoupled}

In Actor-Critic training, the Critic and Actor appear in the same loop, but they are \textbf{two independent optimisation processes}:

\[
\begin{aligned}
  \mathcal{L}_{\text{critic}} &= \bigl(V_\phi(s) - V_{\text{target}}\bigr)^2,
  \qquad &\text{updates only } \phi; \\
  \mathcal{L}_{\text{actor}}  &= -\log\pi_\theta(a|s)\cdot A(s,a),
  \qquad &\text{updates only } \theta.
\end{aligned}
\]

\textbf{Sampling level: coupled.}
Policy determines what trajectories are collected; those trajectories are the data for both Critic and Actor.
Change the policy, and the data distribution changes for everyone.

\textbf{Gradient level: decoupled.}
When computing $\nabla_\theta \mathcal{L}_{\text{actor}}$, the advantage $A(s,a)$ is treated as a \emph{constant}---a \texttt{stop\_gradient} wall separates the two.
The Critic's parameters receive no gradient from the Actor's loss, and vice versa.

\[
  \nabla_\theta J = \mathbb{E}\!\bigl[\nabla_\theta \log\pi_\theta(a|s)\cdot \underbrace{A(s,a)}_{\text{detached}}\bigr]
\]

The one exception is DDPG/TD3's ``max-$Q$ loss'' $\mathcal{L}_{\text{actor}} = -Q(s,\mu_\theta(s))$, where the gradient flows \emph{through} $Q$ to the action and then to the actor---but $Q$'s own parameters are still frozen during the actor update.

This alternating-iteration structure---Critic chases the policy's changing data, Actor chases the Critic's changing landscape---is the root cause of Actor-Critic instability.

%------------------------------------------------------------------
\section{RL Fine-Tuning of SL Base Models}
\label{sec:rl-finetune}

The dominant paradigm for LLM alignment: pretrain with SL, then fine-tune with RL.

\[
  \text{Pretraining (SL)} \;\to\; \text{SFT (SL)} \;\to\; \text{Preference Fine-Tuning (RL)}
\]

\subsection{PPO Approach (InstructGPT Route)}

Requires \textbf{four models} in memory simultaneously:
\begin{enumerate}
  \item \textbf{Policy Model} --- the LLM being trained.
  \item \textbf{Reference Model} --- frozen SFT copy, for KL constraint.
  \item \textbf{Reward Model} --- trained from preferences, frozen as scorer.
  \item \textbf{Critic / Value Head} --- estimates $V(s)$ for advantage computation; typically initialised from the RM.
\end{enumerate}

Training loop: Policy generates $\to$ RM scores $\to$ Critic estimates baseline $\to$ compute advantage $\to$ PPO clipped-ratio update $\to$ KL penalty keeps Policy near Reference.

\subsection{GRPO Approach (DeepSeek Route)}

Eliminates the Critic entirely.
Requires only \textbf{two--three models}: Policy, Reference, and a reward source (RM or rule-based verifier).

Core mechanism:
\begin{enumerate}
  \item For each prompt, Policy generates a \emph{group} of $G$ responses.
  \item Each response is scored.
  \item Within-group relative ranking serves as baseline (no Critic needed).
  \item Policy gradient updates the Policy.
\end{enumerate}

DeepSeek-R1 uses full fine-tuning with GRPO, training all parameters.

\subsection{DPO Approach (Stanford Route)}

Eliminates even the Reward Model.
Requires only \textbf{two models}: Policy and Reference.
Trains directly on preference pairs $(\text{chosen}, \text{rejected})$ with a classification-style loss:
\[
  \mathcal{L}_{\text{DPO}} = -\log\sigma\!\Bigl(\beta\Bigl[
    \log\frac{\pi_\theta(\text{chosen})}{\pi_{\text{ref}}(\text{chosen})}
    - \log\frac{\pi_\theta(\text{rejected})}{\pi_{\text{ref}}(\text{rejected})}
  \Bigr]\Bigr).
\]

\subsection{Comparison}

\begin{table}[htbp]
\centering
\caption{LLM RL fine-tuning approaches.}
\begin{tabular}{@{}lcccc@{}}
\toprule
\textbf{Approach} & \textbf{Models} & \textbf{Bellman eq.} & \textbf{Reward source} & \textbf{Complexity} \\
\midrule
PPO & 4 & Used (Critic) & RM & High \\
GRPO & 2--3 & Not used & RM / rules & Medium \\
DPO & 2 & Not used & Pref.\ pairs & Low \\
\bottomrule
\end{tabular}
\end{table}

\subsection{Parameter-Efficient Updates: LoRA and QLoRA}

Regardless of the RL method, one can choose \emph{how} to update the Policy:

\begin{itemize}
  \item \textbf{Full fine-tuning}: update all parameters. Maximum expressiveness; huge memory.
  \item \textbf{LoRA}: freeze the base, train low-rank matrices $\Delta W = AB$. The Reference Model is free---a forward pass without the adapter \emph{is} the reference.
  \item \textbf{QLoRA}: 4-bit quantised base + LoRA. Enables RLHF on consumer GPUs.
\end{itemize}

Typical choices: PPO with full fine-tuning (InstructGPT) or LoRA; GRPO with full fine-tuning (DeepSeek-R1); DPO + LoRA as the most memory-efficient alignment path.

%------------------------------------------------------------------
\section{RL with Diffusion and Flow Matching Models}
\label{sec:diffusion-rl}

When a diffusion or flow-matching model is used as a \emph{policy}, the multi-step denoising/flow process can be cast as an MDP and fine-tuned with RL.

\subsection{Denoising Chain as MDP}

\begin{keyidea}[DDPO: Denoising Diffusion Policy Optimisation]
Model the denoising chain as an episodic MDP:
\emph{state} = current noisy sample $x_t$;
\emph{action} = one denoising step's output;
\emph{episode} = complete chain $x_T \to x_0$;
\emph{reward} = given only on the final $x_0$.
Apply REINFORCE or PPO to estimate the policy gradient for each step without backpropagating through the entire chain.
\end{keyidea}

DDPO (Black et al., 2023) successfully fine-tuned text-to-image diffusion models with aesthetic scores and human preferences.

\textbf{DPPO} (Ren et al., 2025) adapts PPO specifically for diffusion policies in continuous control, addressing cross-step clipping and high-dimensional action spaces.

\textbf{DRaFT} (Clark et al., 2023) backpropagates reward gradients directly through a truncated denoising chain---more sample-efficient but requires a differentiable RM.

\subsection{Flow Matching and RL}

Flow Matching replaces stochastic SDEs with deterministic ODEs, enabling 5--10$\times$ faster inference.
For RL, each discretised ODE step is treated as an action in the MDP, mirroring the DDPO formulation.

\subsection{Key Challenges}

\begin{table}[htbp]
\centering
\caption{Challenges in RL + diffusion/flow.}
\begin{tabular}{@{}ll@{}}
\toprule
\textbf{Challenge} & \textbf{Common solution} \\
\midrule
Cross-step credit assignment & Episodic MDP formulation \\
High REINFORCE variance & Fewer steps / baselines / truncated backprop \\
Reward differentiability & Policy gradient bypasses this \\
KL regularisation & Penalty similar to RLHF \\
Action-chunk consistency & Overlapping window design \\
\bottomrule
\end{tabular}
\end{table}

%------------------------------------------------------------------
\section{Agent RL vs.\ One-Round RL}
\label{sec:agent-rl}

\subsection{Traditional RL: Single-Round Sequential Interaction}

A continuous time-series rollout:
$s_0 \to a_0 \to r_0 \to s_1 \to a_1 \to r_1 \to \cdots \to s_T$.
Each action directly produces the next state and reward.
Typical scenarios: autonomous driving, robot control, Atari games.

\subsection{Agentic RL: Multi-Round Conversational Interaction}

An episode is a multi-round dialogue:
\begin{enumerate}
  \item Agent observes $\to$ thinks $\to$ calls tools / outputs text $\to$ environment feedback.
  \item Agent observes feedback $\to$ thinks $\to$ acts again $\to$ feedback.
  \item[$\vdots$]
  \item Agent produces final output $\to$ episode ends.
\end{enumerate}
Typical scenarios: coding agents, dialogue systems, tool-using agents.

\subsection{The Fundamental Connection}

\begin{keyidea}[Autonomous Driving Is Agentic RL]
In autonomous driving, each action chunk (a control sequence spanning $\sim$100\,ms--2\,s) is analogous to one ``round'' in an agentic AI system.
After each chunk, the agent observes the new environment state and decides again.
Both are closed-loop MDPs; the differences are only in action granularity, environment nature, and interaction frequency.
\end{keyidea}

\begin{table}[htbp]
\centering
\caption{Traditional RL vs.\ Agentic RL.}
\begin{tabular}{@{}lll@{}}
\toprule
\textbf{Dimension} & \textbf{Autonomous Driving} & \textbf{Agentic AI} \\
\midrule
One ``round'' & Action chunk (100\,ms--2\,s) & Agent output (tool call / text) \\
Env.\ feedback & New sensor observations & Tool results / user replies \\
Episode & One drive (min--hrs) & One task (code fix / Q\&A) \\
State space & Continuous, high-dim & Discrete, variable-length \\
Action space & Continuous & Discrete / mixed \\
\bottomrule
\end{tabular}
\end{table}

\subsection{Training Differences}

\textbf{Traditional RL} (driving / robotics): continuous-action policies (Gaussian / Diffusion / Flow), Actor-Critic (SAC / TD3 / PPO), high-frequency simulator interaction, engineered rewards.

\textbf{Agentic RL} (LLM agents): autoregressive token policy, PPO / GRPO / DPO, multi-round agent--environment interaction, task-completion / preference / verifier rewards, far higher compute per ``action.''

%------------------------------------------------------------------
\section{Sim2Real / Real2Sim Loop}
\label{sec:sim2real}

RL needs massive interaction, but real-world interaction is expensive, dangerous, and slow.
The Sim2Real and Real2Sim loop addresses this.

\subsection{Sim2Real: From Simulation to Reality}

The core challenge is the \textbf{reality gap}---a policy trained in simulation often fails in the real world because physics, rendering, and dynamics differ.

\textbf{Domain Randomisation.}
Randomise simulator parameters (friction, mass, lighting, textures, sensor noise) during training so the real world is ``just another draw'' from the training distribution.
OpenAI's Automatic Domain Randomisation (ADR, 2019) trained a robotic hand to solve a Rubik's cube by dynamically expanding randomisation across thousands of parameters.

\textbf{System Identification.}
Measure real-world physics parameters, calibrate the simulator, and optionally randomise \emph{around} calibrated values rather than uniformly.

\textbf{Domain Adaptation.}
Use CycleGAN or adversarial training to align visual distributions between sim and real without paired data.

\textbf{Progressive Transfer.}
Train in simulation first, then fine-tune in the real environment with a small amount of real data.

\subsection{Real2Sim: From Reality to Simulation}

The reverse direction: use real-world data to build or improve simulators.

\begin{itemize}
  \item \textbf{Neural Scene Reconstruction}: NeRF / 3D Gaussian Splatting reconstruct photo-realistic, editable scenes from real sensor data.
  \item \textbf{Learned World Models}: generative models trained on real data serve as implicit simulators (GAIA-1, DayDreamer).
  \item \textbf{Digital Twins}: high-fidelity digital replicas continuously updated with sensor data.
  \item \textbf{Sim Parameter Calibration}: optimise simulator parameters (friction, drag, latency) using real trajectory data via Bayesian optimisation or differentiable simulation.
\end{itemize}

\subsection{The Continuous Loop}

In autonomous driving, the loop runs continuously:
\[
  \text{Real data} \xrightarrow{\text{Real2Sim}} \text{Better sim}
  \xrightarrow{\text{Train}} \text{Better policy}
  \xrightarrow{\text{Sim2Real}} \text{Deploy}
  \xrightarrow{\text{Collect}} \text{New real data}
  \to \cdots
\]
Each iteration reduces the reality gap and improves real-world performance.

%------------------------------------------------------------------
\section{Summary}

\begin{table}[htbp]
\centering
\caption{RL's role across different scenarios.}
\begin{tabular}{@{}lll@{}}
\toprule
\textbf{Scenario} & \textbf{RL's role} & \textbf{Key methods} \\
\midrule
From scratch & Full policy learning & DQN, PPO, SAC \\
LLM alignment & Preference fine-tuning & PPO+RM, GRPO, DPO (+LoRA) \\
Diffusion / Flow & Denoising-chain MDP & DDPO, DPPO, DRaFT \\
Agent RL & Multi-round closed-loop & PPO, GRPO for LLM agents \\
Sim $\leftrightarrow$ Real & Transfer loop & Domain rand., world models \\
\bottomrule
\end{tabular}
\end{table}

The common thread: RL is \emph{preference optimisation under biased sampling}.
Whether the policy is a Q-network, an LLM, or a diffusion model, the same closed-loop challenge applies---the algorithm changes its own data distribution.
Every technique in this chapter is a different strategy for managing that fundamental difficulty in a different application context.


\part{The Code}
\chapter{Reading Algorithms Through the Lens}
\label{ch:code-reading}

This chapter walks through thirteen algorithms in the order they appeared historically.
For each, we answer four questions:
\begin{enumerate}[label=(\alph*)]
  \item Where does it sit in the macro evolution arc (Part~I)?
  \item Which micro mechanism from Part~II does it embody?
  \item What is the key code pattern (pseudocode / equations)?
  \item What is the single ``aha'' insight?
\end{enumerate}

%==================================================================
\section{DQN --- Deep Q-Network}
\label{sec:dqn}

\paragraph{(a) Macro arc.}
DQN is the bridge from tabular RL to deep RL---the second relaxation in the DP$\to$RL chain (function approximation replaces the Q-table).

\paragraph{(b) Micro mechanisms.}
\emph{Self-bootstrapping MSE}: the model creates its own labels via the Bellman target $y = r + \gamma\max_{a'}Q_{\text{target}}(s',a')$.
\emph{Replay buffer} breaks temporal correlation.
\emph{Target network} stabilises the moving target.

\paragraph{(c) Key code pattern.}

\begin{lstlisting}[language=Python]
# DQN update
batch = replay_buffer.sample(N)
with no_grad():
    q_target = batch.r + gamma * Q_target(batch.s_next).max(dim=-1)
loss = mse_loss(Q(batch.s, batch.a), q_target)
loss.backward(); optimizer.step()
# periodically: Q_target.load_state_dict(Q.state_dict())
\end{lstlisting}

\paragraph{(d) Aha.}
The label is computed by the model itself---this is not ordinary regression but \emph{self-bootstrapping regression}, which is why a target network is indispensable.

%==================================================================
\section{Double DQN}
\label{sec:ddqn}

\paragraph{(a) Macro arc.}
A direct patch to DQN's overestimation problem---still within the ``stabilise greedy'' phase of evolution.

\paragraph{(b) Micro mechanisms.}
\emph{Separation of selection and evaluation}: the online network selects the best action; the target network evaluates it.
This breaks the positive feedback loop where a single noisy $Q$ both picks and scores.

\paragraph{(c) Key code pattern.}

\begin{lstlisting}[language=Python]
# Double DQN target
a_star = Q_online(s_next).argmax(dim=-1)        # select with online
q_target = r + gamma * Q_target(s_next, a_star)  # evaluate with target
loss = mse_loss(Q_online(s, a), q_target)
\end{lstlisting}

\paragraph{(d) Aha.}
$\mathbb{E}[\max \hat{Q}] > \max \mathbb{E}[\hat{Q}]$---taking the max of noisy estimates is \emph{always} biased upward.
Splitting selection from evaluation does not eliminate the bias but breaks the self-reinforcing cycle.

%==================================================================
\section{Dueling DQN}
\label{sec:dueling}

\paragraph{(a) Macro arc.}
An architectural refinement within the value-based family, improving representational efficiency without changing the loss.

\paragraph{(b) Micro mechanisms.}
\emph{Compressed representation}: decompose $Q$ into a state-value stream $V(s)$ and an advantage stream $A(s,a)$:
\[
  Q(s,a) = V(s) + A(s,a) - \frac{1}{|\mathcal{A}|}\sum_{a'} A(s,a').
\]
The subtraction of the mean advantage ensures identifiability.

\paragraph{(c) Key code pattern.}

\begin{lstlisting}[language=Python]
# Dueling architecture forward pass
features = backbone(s)
v = value_head(features)          # shape: (B, 1)
a = advantage_head(features)      # shape: (B, |A|)
q = v + a - a.mean(dim=-1, keepdim=True)
\end{lstlisting}

\paragraph{(d) Aha.}
Many states have similar value regardless of the action taken.
Factoring out $V(s)$ lets the network say ``this state is good'' without having to re-derive that fact separately for every action.

%==================================================================
\section{REINFORCE}
\label{sec:reinforce}

\paragraph{(a) Macro arc.}
The first policy-gradient algorithm---the pivot from value-first to policy-first in the macro evolution.

\paragraph{(b) Micro mechanisms.}
\emph{Weighted log-likelihood}: increase the probability of actions that led to high returns.
The gradient $\nabla_\theta J = \mathbb{E}[\nabla_\theta\log\pi_\theta(a|s)\cdot G_t]$ is derived from the likelihood ratio trick, and the loss is constructed \emph{after the fact} for autograd.

\paragraph{(c) Key code pattern.}

\begin{lstlisting}[language=Python]
# REINFORCE
returns = compute_discounted_returns(rewards)   # G_t
log_probs = [log_pi(a_t | s_t) for (s_t, a_t) in trajectory]
loss = -sum(lp * G for lp, G in zip(log_probs, returns))
loss.backward(); optimizer.step()
\end{lstlisting}

\paragraph{(d) Aha.}
The loss is not the starting point---the \emph{gradient} is.
The loss $-\log\pi\cdot G$ is reverse-engineered so that autograd reproduces the policy gradient theorem.

%==================================================================
\section{A2C --- Advantage Actor-Critic}
\label{sec:a2c}

\paragraph{(a) Macro arc.}
Introduces a learned baseline (Critic) to REINFORCE, reducing variance while staying on-policy.

\paragraph{(b) Micro mechanisms.}
\emph{Actor-Critic split}: the Critic estimates $V(s)$; the advantage $A = r + \gamma V(s') - V(s)$ replaces the raw return $G_t$ as the weight.
Two independent losses, two independent parameter sets, sharing the same batch of on-policy data.

\paragraph{(c) Key code pattern.}

\begin{lstlisting}[language=Python]
# A2C
v = critic(s); v_next = critic(s_next)
advantage = r + gamma * v_next.detach() - v.detach()  # stop_gradient
loss_actor  = -(log_pi(a|s) * advantage).mean()
loss_critic = mse_loss(v, (r + gamma * v_next).detach())
(loss_actor + loss_critic).backward(); optimizer.step()
\end{lstlisting}

\paragraph{(d) Aha.}
The advantage is \texttt{detach()}'d when computing the actor loss---this is the concrete code manifestation of ``coupled in sampling, decoupled in gradients.''

%==================================================================
\section{DDPG --- Deep Deterministic Policy Gradient}
\label{sec:ddpg}

\paragraph{(a) Macro arc.}
Extends Q-learning to continuous action spaces by replacing $\arg\max$ with a learned actor---the ``need an actor'' relaxation.

\paragraph{(b) Micro mechanisms.}
\emph{Max-$Q$ loss}: the actor's loss is $\mathcal{L} = -Q(s,\mu_\theta(s))$, directly maximising the critic's score.
This turns discrete $\arg\max$ into a differentiable gradient-ascent problem.
Off-policy with replay buffer.

\paragraph{(c) Key code pattern.}

\begin{lstlisting}[language=Python]
# DDPG
# Critic update
a_next = actor_target(s_next)
q_target = r + gamma * critic_target(s_next, a_next)
loss_critic = mse_loss(critic(s, a), q_target.detach())

# Actor update
loss_actor = -critic(s, actor(s)).mean()   # gradient flows through Q to actor
loss_actor.backward(); actor_optimizer.step()
\end{lstlisting}

\paragraph{(d) Aha.}
The actor gradient flows \emph{through} the critic: $\nabla_\theta\mu_\theta(s)$ receives $\nabla_a Q(s,a)\big|_{a=\mu_\theta(s)}$.
The critic's \emph{parameters} are frozen during this step, but its \emph{output} is not detached---this is the key difference from A2C's \texttt{stop\_gradient}.

%==================================================================
\section{TD3 --- Twin Delayed DDPG}
\label{sec:td3}

\paragraph{(a) Macro arc.}
Patches DDPG's three failure modes: overestimation, actor exploiting critic errors, and sensitivity to sharp $Q$ peaks.

\paragraph{(b) Micro mechanisms.}
Three stabilisation tricks:
(i) \emph{Clipped double $Q$}: $y = r + \gamma\min(Q_1, Q_2)$ suppresses overestimation.
(ii) \emph{Delayed policy update}: update actor every $d$ critic steps so the critic stabilises first.
(iii) \emph{Target policy smoothing}: add noise to the target action to prevent the actor from exploiting narrow $Q$ peaks.

\paragraph{(c) Key code pattern.}

\begin{lstlisting}[language=Python]
# TD3 critic target
a_next = actor_target(s_next) + clip(noise, -c, c)  # smoothing
q1_tgt = critic1_target(s_next, a_next)
q2_tgt = critic2_target(s_next, a_next)
q_target = r + gamma * min(q1_tgt, q2_tgt)

# Actor update (every d steps only)
if step % d == 0:
    loss_actor = -critic1(s, actor(s)).mean()
\end{lstlisting}

\paragraph{(d) Aha.}
TD3 is a case study in \emph{bias management}: every trick trades a small conservative bias for a large reduction in destructive overestimation.

%==================================================================
\section{TRPO --- Trust Region Policy Optimisation}
\label{sec:trpo}

\paragraph{(a) Macro arc.}
The theoretical foundation for constrained policy updates---ensures monotonic improvement by bounding the KL divergence between old and new policies.

\paragraph{(b) Micro mechanisms.}
\emph{Trust region}: maximise a surrogate objective subject to a KL constraint:
\[
  \max_\theta\; \mathbb{E}\!\Bigl[\frac{\pi_\theta(a|s)}{\pi_{\text{old}}(a|s)} A^{\pi_{\text{old}}}(s,a)\Bigr]
  \quad\text{s.t.}\quad \mathbb{E}\!\bigl[D_{\text{KL}}(\pi_{\text{old}}\|\pi_\theta)\bigr] \le \delta.
\]
Solved via conjugate gradient + line search on the natural gradient.

\paragraph{(c) Key code pattern.}

\begin{lstlisting}[language=Python]
# TRPO (simplified)
ratio = pi_new(a|s) / pi_old(a|s)
surrogate = (ratio * advantage).mean()
g = grad(surrogate, theta)                # policy gradient
Hv = lambda v: fisher_vector_product(v)   # F * v
step_dir = conjugate_gradient(Hv, g)      # F^{-1} g
step_size = sqrt(2*delta / (step_dir @ Hv(step_dir)))
theta += step_size * step_dir             # with line search for KL
\end{lstlisting}

\paragraph{(d) Aha.}
TRPO shows that \emph{how far} you step matters as much as \emph{which direction} you step.
The KL constraint turns an unconstrained, unstable optimisation into a monotonically improving sequence.

%==================================================================
\section{PPO --- Proximal Policy Optimisation}
\label{sec:ppo}

\paragraph{(a) Macro arc.}
The practical successor to TRPO---replaces the expensive conjugate-gradient solve with a simple clipped surrogate, making trust-region ideas accessible to large-scale training.

\paragraph{(b) Micro mechanisms.}
\emph{Clipped ratio}: the importance ratio $r(\theta) = \pi_\theta(a|s)/\pi_{\text{old}}(a|s)$ is clipped to $[1{-}\varepsilon,\,1{+}\varepsilon]$, preventing any single sample from causing a destructively large update.
Near on-policy: data from $\pi_{\text{old}}$, reused for a few epochs.

\paragraph{(c) Key code pattern.}

\begin{lstlisting}[language=Python]
# PPO clipped objective
ratio = pi_theta(a|s) / pi_old(a|s)
clip_ratio = clip(ratio, 1-eps, 1+eps)
loss = -min(ratio * A, clip_ratio * A).mean()
loss.backward(); optimizer.step()
\end{lstlisting}

\paragraph{(d) Aha.}
PPO deliberately introduces bias (clipping discards some gradient information) to buy stability.
It is the canonical example of \emph{trading bias for engineering reliability}.

%==================================================================
\section{SAC --- Soft Actor-Critic}
\label{sec:sac}

\paragraph{(a) Macro arc.}
The culmination of the ``abandon hard argmax'' trend---entropy becomes part of the objective, not an external trick.

\paragraph{(b) Micro mechanisms.}
\emph{Maximum entropy RL}: the objective is
$J = \mathbb{E}\bigl[\sum_t r_t + \alpha\,\mathcal{H}[\pi(\cdot|s_t)]\bigr]$.
Clipped double $Q$ (from TD3) for the critic.
Automatic temperature $\alpha$ tuning to maintain a target entropy.
Off-policy with replay buffer.

\paragraph{(c) Key code pattern.}

\begin{lstlisting}[language=Python]
# SAC actor update
a_new, log_prob = pi.sample(s)   # reparameterised
q1, q2 = critic1(s, a_new), critic2(s, a_new)
loss_actor = (alpha * log_prob - min(q1, q2)).mean()

# SAC critic update
with no_grad():
    a_next, log_prob_next = pi.sample(s_next)
    q_target = r + gamma * (min(Q1_tgt(s_next, a_next),
                                Q2_tgt(s_next, a_next))
                            - alpha * log_prob_next)
loss_critic = mse_loss(Q(s, a), q_target)
\end{lstlisting}

\paragraph{(d) Aha.}
SAC's entropy-augmented objective is a \emph{biased} objective on purpose.
It does not maximise raw reward but entropy-regularised reward---and this intentional bias makes exploration natural, $Q$-learning smoother, and the policy resistant to premature collapse.

%==================================================================
\section{GRPO --- Group Relative Policy Optimisation}
\label{sec:grpo}

\paragraph{(a) Macro arc.}
A simplification of PPO for LLM alignment that eliminates the Critic, moving further along the ``fewer models'' trend from PPO $\to$ GRPO $\to$ DPO.

\paragraph{(b) Micro mechanisms.}
\emph{Group-relative baseline}: instead of a learned $V(s)$, GRPO generates $G$ completions per prompt and uses within-group reward statistics (mean, std) as the baseline.
This turns the advantage into a \emph{relative ranking} within the group:
\[
  A_i = \frac{R_i - \text{mean}(\{R_j\}_{j=1}^G)}{\text{std}(\{R_j\}_{j=1}^G)}.
\]
KL penalty against a reference model prevents drift.

\paragraph{(c) Key code pattern.}

\begin{lstlisting}[language=Python]
# GRPO
responses = [policy.generate(prompt) for _ in range(G)]
rewards = [reward_fn(prompt, resp) for resp in responses]
advantages = (rewards - mean(rewards)) / (std(rewards) + eps)
for resp, A in zip(responses, advantages):
    ratio = pi_theta(resp|prompt) / pi_old(resp|prompt)
    loss = -clip(ratio, 1-eps, 1+eps) * A
    loss += beta * kl_divergence(pi_theta, pi_ref)
loss.backward(); optimizer.step()
\end{lstlisting}

\paragraph{(d) Aha.}
You do not need a learned value function to compute advantage---comparing siblings from the same prompt provides a surprisingly effective baseline, because it automatically controls for prompt difficulty.

%==================================================================
\section{DPO --- Direct Preference Optimisation}
\label{sec:dpo}

\paragraph{(a) Macro arc.}
The furthest point on the simplification axis: no RM, no Critic, no RL loop---preference learning reduced to a single classification loss.

\paragraph{(b) Micro mechanisms.}
\emph{Implicit reward}: DPO derives that the optimal RL policy under a Bradley--Terry preference model satisfies
$r(x,y) = \beta\log\frac{\pi(y|x)}{\pi_{\text{ref}}(y|x)} + C$.
Substituting this into the preference loss yields a closed-form objective that depends only on the policy and reference:
\[
  \mathcal{L}_{\text{DPO}} = -\log\sigma\!\Bigl(\beta\Bigl[
    \log\frac{\pi_\theta(y_w|x)}{\pi_{\text{ref}}(y_w|x)}
    - \log\frac{\pi_\theta(y_l|x)}{\pi_{\text{ref}}(y_l|x)}
  \Bigr]\Bigr).
\]

\paragraph{(c) Key code pattern.}

\begin{lstlisting}[language=Python]
# DPO
logr_w = log_pi_theta(y_win|x) - log_pi_ref(y_win|x)
logr_l = log_pi_theta(y_lose|x) - log_pi_ref(y_lose|x)
loss = -log_sigmoid(beta * (logr_w - logr_l))
loss.backward(); optimizer.step()
\end{lstlisting}

\paragraph{(d) Aha.}
DPO shows that the reward model is not a separate entity---it is \emph{implicit} in the ratio $\pi/\pi_{\text{ref}}$.
By algebraically eliminating the RM, DPO collapses the entire RLHF pipeline into a single supervised loss on preference pairs.

%==================================================================
\section{ProRL --- Process-Reward RL}
\label{sec:prorl}

\paragraph{(a) Macro arc.}
ProRL represents the frontier of RL for reasoning: instead of scoring only the final answer (outcome reward), it provides reward at each reasoning step (process reward), enabling finer-grained credit assignment in chain-of-thought settings.

\paragraph{(b) Micro mechanisms.}
\emph{Process Reward Model (PRM)}: a learned model that scores each intermediate step of a reasoning chain, trained on step-level annotations.
This converts a sparse-reward problem (correct/incorrect final answer) into a dense-reward MDP where each reasoning step receives feedback.
Combined with policy gradient (PPO or GRPO variants), the per-step reward dramatically improves credit assignment.

\paragraph{(c) Key code pattern.}

\begin{lstlisting}[language=Python]
# ProRL with Process Reward Model
steps = policy.generate_chain_of_thought(prompt)  # list of steps
step_rewards = [PRM.score(prompt, steps[:i+1]) for i in range(len(steps))]
# Compute per-step advantage using step-level rewards
advantages = compute_gae(step_rewards, values, gamma, lam)
for step, A in zip(steps, advantages):
    ratio = pi_theta(step|context) / pi_old(step|context)
    loss = -clip(ratio, 1-eps, 1+eps) * A
loss.backward(); optimizer.step()
\end{lstlisting}

\paragraph{(d) Aha.}
Outcome reward tells you \emph{whether} you got it right; process reward tells you \emph{where} you went wrong.
By rewarding each reasoning step, ProRL turns the credit-assignment problem in long chains of thought from a needle-in-a-haystack search into a step-by-step guided climb.

%==================================================================
\section{Summary: The Evolution at a Glance}

\begin{table}[htbp]
\centering
\caption{Thirteen algorithms mapped to the evolution arc and micro mechanisms.}
\small
\begin{tabular}{@{}rlll@{}}
\toprule
\textbf{\#} & \textbf{Algorithm} & \textbf{Core micro mechanism} & \textbf{Key trade-off} \\
\midrule
1  & DQN        & Self-bootstrapping MSE + replay            & Bias from moving target \\
2  & Double DQN & Separate selection / evaluation             & Reduced overestimation \\
3  & Dueling DQN & $V$/$A$ decomposition                     & Better state-value sharing \\
4  & REINFORCE  & Weighted log-likelihood                     & High variance \\
5  & A2C        & Critic as baseline                          & Two-optimizer instability \\
6  & DDPG       & Max-$Q$ loss for continuous actions          & Actor exploits critic \\
7  & TD3        & Twin critics + delayed update               & Conservative bias \\
8  & TRPO       & KL-constrained trust region                 & Expensive second-order solve \\
9  & PPO        & Clipped importance ratio                    & Bias for stability \\
10 & SAC        & Max-entropy objective                       & Entropy-biased reward \\
11 & GRPO       & Group-relative baseline                     & No learned value function \\
12 & DPO        & Implicit reward in $\pi/\pi_{\text{ref}}$  & No online RL loop \\
13 & ProRL      & Step-level process reward                   & Needs step-level annotations \\
\bottomrule
\end{tabular}
\end{table}

Reading these thirteen algorithms in sequence reveals a clear narrative: the field moved from \emph{learning values and extracting behaviour via argmax} toward \emph{directly shaping preference distributions}, while progressively simplifying the training infrastructure from four models (PPO) to two (DPO) to dense per-step feedback (ProRL).
Every ``trick''---replay, target networks, double Q, clipping, entropy, group baselines---is a specific response to a specific failure mode in the closed loop where the policy shapes its own training data.
Once you see this pattern, no RL algorithm is truly mysterious.


% ============================================================
\backmatter
% ============================================================

\chapter*{Notation}
\addcontentsline{toc}{chapter}{Notation}

\begin{tabular}{ll}
\toprule
\textbf{Symbol} & \textbf{Meaning} \\
\midrule
$s, s'$ & State, next state \\
$a, a'$ & Action, next action \\
$r$ & Reward \\
$\gamma$ & Discount factor \\
$\pi(a|s)$ & Stochastic policy \\
$\mu(s)$ & Deterministic policy \\
$\pi_\theta$ & Policy parameterized by $\theta$ \\
$V^{\pi}(s)$ & State-value function under $\pi$ \\
$Q^{\pi}(s,a)$ & Action-value function under $\pi$ \\
$V^{\ast}(s), Q^{\ast}(s,a)$ & Optimal value functions \\
$A^{\pi}(s,a)$ & Advantage function: $Q^{\pi}(s,a) - V^{\pi}(s)$ \\
$G_t$ & Discounted return from time $t$ \\
$\mathbb{E}[\cdot]$ & Expectation \\
$\hat{\mathbb{E}}[\cdot]$ & Sample estimate of expectation \\
$\mathcal{H}(\pi)$ & Entropy of policy $\pi$ \\
$D_{\mathrm{KL}}$ & Kullback-Leibler divergence \\
$\alpha$ & Learning rate or temperature parameter (SAC) \\
$\epsilon$ & Exploration rate ($\epsilon$-greedy) or clipping range (PPO) \\
\bottomrule
\end{tabular}

\end{document}
